% !TEX program = xelatex
\documentclass[UTF8,a4paper,11pt,fontset=windows]{ctexart}
\usepackage[margin=2.2cm]{geometry}
\usepackage{amsmath,amssymb,amsthm,enumitem,booktabs,longtable}
\usepackage[unicode,hidelinks]{hyperref}
\usepackage{xurl}
\usepackage{fancyhdr}
\setlength{\headheight}{15pt}
\pagestyle{fancy}\fancyhf{}\fancyhead[L]{抽象代数集中测验复习}\fancyhead[R]{往次作业详解}\fancyfoot[C]{\thepage}
\setlength{\parindent}{2em}\setlength{\parskip}{0.35em}
\setlist[enumerate]{itemsep=0.35em,topsep=0.4em}
\renewcommand{\proofname}{证明}
\theoremstyle{definition}
\newtheorem{reviewtheorem}{定理}
\newtheorem{reviewprop}[reviewtheorem]{命题}
\newtheorem{reviewlemma}[reviewtheorem]{引理}
\newtheorem{reviewdefinition}[reviewtheorem]{定义}
\renewcommand{\thereviewtheorem}{K\arabic{reviewtheorem}}
\newcommand{\toolref}[1]{\hyperref[tool:#1]{结论~\ref*{tool:#1}}}
\newcommand{\Z}{\mathbb Z}\newcommand{\Q}{\mathbb Q}
\newcommand{\R}{\mathbb R}\newcommand{\C}{\mathbb C}
\newcommand{\N}{\mathbb N}\newcommand{\F}{\mathbb F}
\DeclareMathOperator{\GL}{GL}\DeclareMathOperator{\SL}{SL}
\DeclareMathOperator{\Aut}{Aut}\DeclareMathOperator{\Inn}{Inn}
\DeclareMathOperator{\lcm}{lcm}\DeclareMathOperator{\ord}{ord}
\DeclareMathOperator{\sgn}{sgn}\DeclareMathOperator{\Ker}{Ker}
\title{抽象代数集中测验复习资料\\\large 往次作业题目、考点与详细证明}
\author{}\date{整理日期：2026年10月26日\quad 测验日期：2026年10月30日（周五）}
\begin{document}
\maketitle
\section*{使用说明与范围核对}
本资料按提供的四张作业截图确定范围，共54题：第零章9题、第一章19题、第二章22题、第三章4题。各章均保留教材习题号，不按本资料中的节号改编号。截图中的“调休补课作业39、43、44”归于其前文第一章；另列的第二章38、43、44归于第二章。

题面以聂灵沼、丁石孙《代数学引论》第三版为准：第零章书页20--22、第一章54--58、第二章98--100、第三章131--132；扫描PDF页码均为书页加10。第一至第三章详解以现有笔记的原解为基础，保留其构造、证明方式和推理顺序，在原思路内补齐步骤；第零章沿用教材的定义、递归构造及相关定理证明方法。每题的“原文依据”说明实际对照的定义、定理或笔记解答。

\textbf{关于补充证明：}上述习题页没有给出这些题目的完整证明。教材正文已有证明或构造的，依照其路线展开；正文省略证明的，沿原定义和已介绍的方法补全，并标为补充证明。以现有笔记原解为依据的，不冒称为教材原证。

\textbf{题面订正：}第二章习题12已据教材恢复：第一个生成元为$(123)$；第二个生成元在奇数情形为$(12\cdots n)$，在偶数情形为$(23\cdots n)$。现有笔记中的另一组生成元并非教材原题，本资料已纠正；详解仍沿用共轭与逐步加入新点的原有思路。

\begin{center}\small
\begin{tabular}{p{1.6cm}p{11.6cm}}
\toprule 章节 & 截图指定题号\\\midrule
第零章 & 4--8、13、14、23、24\\
第一章 & 1、2、4、6、7--10、12、18、21、22、29、30、32、34、39、43、44\\
第二章 & 7、8、10--12、15、29--32、36、38--47、49\\
第三章 & 1、2、3、8\\\bottomrule
\end{tabular}
\end{center}

\textbf{复习顺序：}先独立写出每题的构造或第一步，再核对详细证明。第一章重点练习子群与理想判别、同态构造；第二章重点练习生成元、群作用、西罗计数与可解性；第三章重点练习非交换乘法次序、单侧逆与反例。没有依据把某题标为“必考”，各题均在已给出的作业范围内。

\textbf{记号约定：}群的单位元记为$e$，环的单位元记为$1$；$|G|$为群的阶，$\ord(a)$为元素阶，$\langle S\rangle$为包含$S$的最小子群。$H\leq G$表示子群，$N\trianglelefteq G$表示正规子群，$[G:H]$为左陪集数。置换乘法右边先作用。一般环不预设有单位元；涉及$1$和单位时另作幺环假设。第一章涉及子环时不要求与母环共用单位元。

\tableofcontents\clearpage
\section*{证明工具与引用标识}
\addcontentsline{toc}{section}{证明工具与引用标识}
下列定理、命题和引理采用本复习册独立的K编号；K编号不冒充原笔记或教材编号。每条给出适用条件、完整结论和可回查的来源。各题“所用结论”列出实际依据，点击编号可跳转到对应表述。

工具在这里集中列出以便查阅，不改变习题的证明先后关系。标为“本题所证”或“补充引理”的条目，证明仍须在相应习题中完成；不得将它们反过来当作该题已经成立的前提。环和群的基本定义沿用前面的记号约定及各章说明。
\setcounter{secnumdepth}{0}
% !TEX root = main.tex
\subsection{集合、自然数与整数工具}

\begin{reviewdefinition}[关系的自反、对称与传递]\label{tool:n-relations}
设$S$是集合，$\sim$是$S$上的二元关系。自反性指对所有$x\in S$有$x\sim x$；对称性指对所有$x,y\in S$，$x\sim y$蕴含$y\sim x$；传递性指对所有$x,y,z\in S$，$x\sim y$且$y\sim z$蕴含$x\sim z$。同时满足三者称为等价关系。否定任一性质只需给出违反其全称命题的实例；对称性、传递性的反例须满足相应蕴含式的前提。若前提对所有元素均不成立，则这条蕴含式成立。

\textbf{出处：}教材第零章第1节等价关系定义，书页2；本册第零章习题4、5直接按这三条定义核对。
\end{reviewdefinition}

\begin{reviewtheorem}[数学归纳法]\label{tool:n-induction}
设$P(n)$是关于$n\in\N=\{0,1,2,\ldots\}$的命题。若$P(0)$成立，且对每个$n\in\N$，$P(n)$成立蕴含$P(n+1)$成立，则$P(n)$对所有自然数成立。从$n_0$开始的归纳须检验$P(n_0)$并对全部$n\ge n_0$检验归纳步。强归纳形式是：对每个$n$，若所有$k<n$的$P(k)$成立便能推出$P(n)$，则所有$P(n)$成立。

\textbf{出处：}教材第零章第2节自然数公理与归纳法，书页4--6：第一归纳法见书页4，强归纳法见书页6。该部分尚无对应的原代数笔记章节，本册据已核对教材补入。
\end{reviewtheorem}

\begin{reviewtheorem}[自然数递归构造]\label{tool:n-recursion}
给定集合$X$、元素$a\in X$和映射$\varphi:X\to X$，存在唯一映射$f:\N\to X$使$f(0)=a$、$f(n+1)=\varphi(f(n))$。据此自然数加法、乘法分别由
\[
m+0=m,\quad m+(n+1)=(m+n)+1,\qquad
m\cdot0=0,\quad m(n+1)=mn+m
\]
定义。此条只给递归定义，不预先把加法、乘法的交换律、结合律和分配律作为习题6的前提；这些运算律在该题内按原定义补证。

\textbf{出处：}教材第零章第2节递归定理及运算定义，书页5--6。
\end{reviewtheorem}

\begin{reviewprop}[自然数的次序]\label{tool:n-order}
对$a,b\in\N$，$a\le b$当且仅当存在$t\in\N$使$a+t=b$；$a<b$当且仅当可取$t>0$。任意两自然数满足$a<b$、$a=b$、$b<a$恰有一项成立。次序可传递，且加同一个自然数保持严格不等式。正整数乘积为正可由乘法递归定义归纳证明，习题8中保留了该步，未用乘法消去律反证它。

\textbf{出处：}教材第零章第2节次序与良序性的讨论，书页6；本册第零章习题8补证正整数乘积为正。
\end{reviewprop}

\begin{reviewtheorem}[带余除法与欧几里得算法]\label{tool:n-gcd}
对整数$a$及正整数$b$，存在唯一整数$q,r$使$a=qb+r$、$0\le r<b$。若$a=qb+r$，则$a,b$的公共因子恰为$b,r$的公共因子，因而$\gcd(a,b)=\gcd(b,r)$。重复带余除法，非零余数严格递减，故有限步终止，最后非零余数是正最大公因数。非全为零的有限个整数也有正最大公因数。

\textbf{出处：}教材第零章第3节整除、最大公因数及辗转相除法，书页9--10；现有笔记第二章习题7使用同一整数工具。
\end{reviewtheorem}

\begin{reviewtheorem}[贝祖等式与互素整除]\label{tool:n-bezout}
整数$a_1,\ldots,a_t$不全为零，令$d=\gcd(a_1,\ldots,a_t)>0$，则存在整数$u_i$使$d=\sum_i u_ia_i$；全部整数线性组合构成$d\Z$。特别地，$a,b$互素当且仅当$au+bv=1$有整数解。若$\gcd(a,b)=1$且$a\mid bc$，则$a\mid c$。若两两互素的正整数$m_1,\ldots,m_r$均整除$c$，则$\prod_i m_i\mid c$。

\textbf{出处：}教材第零章第3节定理2、定理3，书页10；多整数形式由两整数形式迭代。现有笔记第二章习题7、第二章循环群一节沿用此性质。
\end{reviewtheorem}

\begin{reviewtheorem}[算术基本定理与素因子]\label{tool:n-prime-factor}
每个大于$1$的整数$n$可唯一写成$n=p_1^{\alpha_1}\cdots p_r^{\alpha_r}$，其中$p_i$是两两不同的素数、$\alpha_i\ge1$，唯一性忽略因子排列。特别地，$n>1$总有素因子；$n$的正约数均且仅为$\prod_i p_i^{\beta_i}$，其中$0\le\beta_i\le\alpha_i$。无平方因子的约数对应$\beta_i\in\{0,1\}$，即不同素因子集合的子集。

\textbf{出处：}教材第零章第3节算术基本定理及标准分解式，书页11--12；本册第零章习题13、14使用。
\end{reviewtheorem}

\begin{reviewprop}[同余运算、逆元及一次方程]\label{tool:n-congruence}
固定正整数$m$，$a\equiv b\pmod m$指$m\mid a-b$。同余可相加、相减、相乘，也可取相同正整数次幂。若$\gcd(a,m)=1$，则$a$在模$m$下有唯一逆元剩余类；若$au+mv=1$，则逆元为$u\pmod m$。此时$ax\equiv b\pmod m$的全部解为$x\equiv ub\pmod m$。同余消去只有在乘数与模数互素时才可直接使用：$ac\equiv ad$且$\gcd(a,m)=1$推出$c\equiv d$。

\textbf{出处：}教材第零章第4节同余性质与定理4，书页12--13；构造路线为贝祖等式后乘以右端常数。
\end{reviewprop}

\begin{reviewtheorem}[孙子定理及原文构造]\label{tool:n-crt}
设$m_1,\ldots,m_r>1$两两互素，$M=\prod_i m_i$。对任意整数$b_i$，方程组$x\equiv b_i\pmod{m_i}$有且仅有一个模$M$的解。令$M_i=M/m_i$，由互素性选择整数$u_i$使$u_iM_i\equiv1\pmod{m_i}$，并令$e_i=u_iM_i$，则
\[
e_i\equiv1\pmod{m_i},\qquad e_i\equiv0\pmod{m_j}\ (j\ne i),\qquad
x\equiv\sum_i b_i e_i\pmod M.
\]
任意两整数解之差被所有$m_i$整除，故被$M$整除，给出唯一性。

\textbf{出处：}教材第零章第4节定理5及其证明，书页14；原文先给两模数构造再用归纳推广，本册习题24保持该构造。
\end{reviewtheorem}

\begin{reviewlemma}[有限集计数与二项式计数]\label{tool:n-binomial-count}
以下集合均有限。分成互不相交的部分时，总基数等于各部分基数之和；逐步选择的每一步若有与此前选择无关的固定选择数$c_i$，总数为$\prod_i c_i$。满射$f:X\to Y$若每个纤维均有$d$个元素，则$|X|=d|Y|$。同一有限集合$S$的子集$A,B$若$|A|+|B|>|S|$，则$A\cap B\ne\varnothing$。多个集合并集的大小不超过各基数之和；若所有集合均含某个共同元素，可将此元素只计一次。

$r$元集合有$\binom rk$个$k$元子集，故在交换数系中
\[
(x+y)^r=\sum_{k=0}^r\binom rkx^ky^{r-k},\qquad
\sum_{k=0}^r\binom rk(-1)^k=(1-1)^r.
\]
\textbf{出处：}前半为本册习题的计数步骤汇总（第一章习题21、44，第二章习题29、30、32），不是另称原书编号定理；末式为第零章习题14使用的二项式计数。环中的无幺元版本另见群与环基础工具。
\end{reviewlemma}

% !TEX root = main.tex
\subsection{群与环的基础工具}
本节完整陈述第一章详解及后续部分所用的基础结论。出处指实际核对的现有笔记；行号是对应源文件的定位，不是原书定理编号。为独立阅读而合并、展开的陈述标明其补充性质。

\begin{reviewprop}[群的消去律、逆元与共轭]\label{tool:g-laws}
设$G$为群，单位元为$e$。单位元和每个元素的逆元均唯一，且对任意$a,b,c\in G$，
\[
(a^{-1})^{-1}=a,\quad (ab)^{-1}=b^{-1}a^{-1},\quad
ab=ac\Rightarrow b=c,\quad ba=ca\Rightarrow b=c.
\]
方程$ax=b$、$ya=b$各有唯一解，分别为$x=a^{-1}b$、$y=ba^{-1}$。
整数幂满足$a^ma^n=a^{m+n}$、$(a^m)^n=a^{mn}$；若$ab=ba$，则$(ab)^m=a^mb^m$对所有整数$m$成立。
固定$c\in G$，共轭映射$x\mapsto c^{-1}xc$是群的自同构；特别地，
\[
c^{-1}(xy)c=(c^{-1}xc)(c^{-1}yc),\qquad
c^{-1}x^mc=(c^{-1}xc)^m.
\]
\textbf{出处：}现有笔记抽象代数第一章第2节“群的定义和简单性质”，单位元与逆元性质、方程唯一解命题、整数指数法则（\path{chapter1.tex:77,90,109}）；共轭保幂的具体应用见本章习题7（\path{chapter1-exercises.tex:59}），此处将该计算完整陈述为共轭保运算结论。
\end{reviewprop}

\begin{reviewtheorem}[子群判别法]\label{tool:g-subgroup}
设$G$为群，$H\subseteq G$非空。以下条件等价：$H$在原运算下为群$G$的子群；对任意$x,y\in H$均有$xy^{-1}\in H$；$H$含$e$且对乘积与取逆封闭。
\textbf{出处：}现有笔记抽象代数第一章第4节“子群、陪集”的子群判别法（\path{chapter1.tex:261}）。
\end{reviewtheorem}

\begin{reviewlemma}[有限满射的等纤维计数]\label{tool:g-fiber-count}
设$X,Y$为有限集合，$f:X\to Y$满射。若存在正整数$d$使每个$y\in Y$的原像集合$f^{-1}(\{y\})$都恰有$d$个元素，则
\[
|X|=d|Y|.
\]
更一般地，$|X|=\sum_{y\in Y}|f^{-1}(\{y\})|$，因为这些原像集合两两不交且其并为$X$。两个有限集合的笛卡尔积满足$|A\times B|=|A||B|$。
\textbf{出处：}现有笔记抽象代数第一章习题21的原像计数（\path{chapter1-exercises.tex:209}）实际使用的有限集合计数原理；本条是对该步骤的补充一般陈述，不标作代数正文已有的独立引理。
\end{reviewlemma}

\begin{reviewprop}[陪集的划分及左右指数]\label{tool:g-cosets}
设$H\leq G$。左陪集$aH=\{ah:h\in H\}$满足
\[
aH=bH\iff a^{-1}b\in H;
\]
右陪集$Ha$满足$Ha=Hb\iff ab^{-1}\in H$。任意两个左陪集相等或不交，右陪集亦然。全部不同左陪集将$G$划分，全部不同右陪集亦将$G$划分；每个陪集与$H$等势。
映射$aH\mapsto Ha^{-1}$给出左陪集集合与右陪集集合的双射，故左右陪集数相同，记为$[G:H]$，即子群指数。这些结论不要求$G$有限。
\textbf{出处：}现有笔记抽象代数第一章第4节“子群、陪集”之“陪集及其计数”（\path{chapter1.tex:309,321,332}）。
\end{reviewprop}

\begin{reviewtheorem}[拉格朗日定理与元素的阶]\label{tool:g-lagrange}
若$G$为有限群、$H\leq G$，则$|G|=[G:H]|H|$，所以$|H|$整除$|G|$。若$a\in G$的阶为$d$，则$d$整除$|G|$，且$a^{|G|}=e$。
一般群中若$a$的阶为有限正整数$d$，则$a^m=e$当且仅当$d\mid m$，并且$a^i=a^j$当且仅当$d\mid(i-j)$。
特别地，素数$p$阶群中任意非单位元的阶都为$p$，所以群为循环群。
\textbf{出处：}现有笔记抽象代数第一章第4节“子群、陪集”的拉格朗日定理、元素阶命题及推论（\path{chapter1.tex:343,359,374}）；素数阶的推论见习题26（\path{chapter1-exercises.tex:263}附近的习题24--26合并解答）。
\end{reviewtheorem}

\begin{reviewprop}[循环子群与无限阶元素]\label{tool:g-cyclic}
对群$G$中的$a$，$\langle a\rangle=\{a^m:m\in\Z\}$是包含$a$的最小子群，且
\[
G=\bigcup_{a\in G}\langle a\rangle.
\]
若$a$无限阶，则$a^i=a^j$蕴含$i=j$，映射$m\mapsto a^m$给出$\Z\cong\langle a\rangle$；各子群$\langle a^m\rangle$（正整数$m$）两两不同。更完整地，无限循环群的任一子群为平凡子群或$\langle a^d\rangle$，其中$d$为唯一正整数；有限阶为$d$的元素则生成恰有$d$个元素的子群。
\textbf{出处：}现有笔记抽象代数第一章第4节的循环子群定义与元素阶命题（\path{chapter1.tex:353,359}）；无限阶产生无限多个子群及全群为循环子群之并见习题18（\path{chapter1-exercises.tex:171}）。无限循环群子群的完整形式是在同一指数思路下补充的陈述，可由子群内最小正指数与整数带余除法得到，未冒称该处原文逐字列出。
\end{reviewprop}

\begin{reviewprop}[正规子群判别及正规子群的乘积]\label{tool:g-normal}
设$H\leq G$。以下条件等价：对任意$g\in G$均有$gH=Hg$；对任意$g$均有$gHg^{-1}=H$；对任意$g$均有$gHg^{-1}\subseteq H$。满足这些条件时称$H$正规，记为$H\trianglelefteq G$。
若$M,N\trianglelefteq G$，则$MN=NM$且$MN\trianglelefteq G$；若$N\trianglelefteq G$且$N\leq H\leq G$，则$N\trianglelefteq H$。
\textbf{出处：}现有笔记抽象代数第一章第6节“同态、正规子群”的正规性等价条件（\path{chapter1.tex:578}）；正规子群乘积见习题22（\path{chapter1-exercises.tex:222}）。将共轭条件限制至$H$得到最后一句，是同一条件的直接应用。
\end{reviewprop}

\begin{reviewtheorem}[群同态基本定理]\label{tool:g-first-iso}
若$\varphi:G\to G'$为群同态，则$\ker\varphi=\{g:\varphi(g)=e'\}$为正规子群，$\varphi(G)\leq G'$，并有良定义的群同构
\[
G/\ker\varphi\longrightarrow\varphi(G),\qquad
g\ker\varphi\longmapsto\varphi(g).
\]
若$\varphi$满射，则$G/\ker\varphi\cong G'$。同态为单射当且仅当其核为$\{e\}$；双射同态就是群同构。
\textbf{出处：}现有笔记抽象代数第一章第6节“同态、正规子群”的核与像性质，以及第7节“商群”的群同态基本定理（\path{chapter1.tex:487,525,578,731}）。
\end{reviewtheorem}

\begin{reviewprop}[子环与理想的判别、单位的作用]\label{tool:g-ring-criteria}
设$R$为环，这里不预设交换或有单位元。非空$S\subseteq R$是子环，当且仅当$a,b\in S$蕴含$a-b\in S$、$ab\in S$。子环若有单位元，不要求它与母环的单位元相同。
非空$I\subseteq R$是双边理想，当且仅当对$a,b\in I$有$a-b\in I$，并且对$r\in R,a\in I$有$ra,ar\in I$。理想对整数倍与有限和封闭，也对左右乘以任意环元封闭。
若$R$有单位元且$1\in I$，或$I$包含$R$的某个单位，则$I=R$。体中任意非零元素可逆，所以体的双边理想只有零理想和全体；域是乘法交换的体。
\textbf{出处：}现有笔记抽象代数第一章第8节“环、子环”的子环判别法（\path{chapter1.tex:888}），第9节“各种特殊类型的环”的体与域，第10节“环的同态、理想”的理想定义与含单位的理想命题（\path{chapter1.tex:1201,1235}）；体无非平凡理想见习题32（\path{chapter1-exercises.tex:332}）。
\end{reviewprop}

\begin{reviewprop}[由双射运输环运算]\label{tool:g-transport}
设$R$为环、$X$为集合，$f:X\to R$为双射。在$X$上定义
\[
x\oplus y=f^{-1}(f(x)+f(y)),\qquad
x\odot y=f^{-1}(f(x)f(y)).
\]
则$X$成为环，$f$是环同构；其零元是$f^{-1}(0)$，$x$的加法逆元是$f^{-1}(-f(x))$。若$R$有单位元，则$X$的单位元是$f^{-1}(1)$；若$R$交换，则$X$也交换。
等价地，若已定义$X$上的运算，验证双射$f$保两种运算即可把这些环公理从$R$运输到$X$。
\textbf{出处：}现有笔记抽象代数第一章习题29的$f(a)=1-a$构造（\path{chapter1-exercises.tex:290}），并结合第8节“环、子环”的环同构讨论（\path{chapter1.tex:930}起）。本条为原构造的完整一般化陈述，不标作原文独立定理。
\end{reviewprop}

\begin{reviewlemma}[交换环二项式展开与理想中的幂]\label{tool:g-binomial}
设$R$为交换环、$I\lhd R$，$r^n\in I$、$s^m\in I$，其中$n,m$为正整数。则
\[
(r\pm s)^{n+m-1}\in I,\qquad
(tr)^n=t^nr^n\in I\quad(t\in R).
\]
所用的二项式展开为
\[
(r+s)^N=\sum_{k=0}^{N}\binom Nk r^ks^{N-k}.
\]
若$R$无单位元，端项分别理解为$s^N,r^N$，不在环内使用零次幂。中间项的整数系数理解为加法整数倍。对$N=n+m-1$，每一项至少含$n$个$r$或$m$个$s$，故由吸收性落在$I$。特别地，若$r^n=s^m=0$，则$(r\pm s)^{n+m-1}=0$。
\textbf{出处：}现有笔记抽象代数第一章习题34“理想的根”和习题39“交换环的幂零元构成理想”（\path{chapter1-exercises.tex:350,396}）；这是将原解实际使用的二项式步骤抽出陈述。
\end{reviewlemma}

\begin{reviewlemma}[同一元素的有限几何级数]\label{tool:g-geometric}
设$R$为幺环，$a\in R$，$n\geq1$。不要求$R$交换，仍有
\[
(1-a)(1+a+\cdots+a^{n-1})
=(1+a+\cdots+a^{n-1})(1-a)=1-a^n.
\]
若$a^n=0$，则$1-a$为单位，其双边逆元为$1+a+\cdots+a^{n-1}$。
\textbf{出处：}现有笔记抽象代数第一章习题38“$1-a$的有限几何级数逆”（\path{chapter1-exercises.tex:388}）。本条供后续环论习题复用。
\end{reviewlemma}

\begin{reviewprop}[有限域中的线性无关列与有序基计数]\label{tool:g-basis-count}
设$F$为有$q$个元素的有限域，$n\geq1$。$F^n$有$q^n$个向量；若$v_1,\ldots,v_k$线性无关，则其张成子空间恰有$q^k$个向量，因为每个向量的$k$个系数唯一。
矩阵可逆当且仅当其列向量构成$F^n$的一组有序基。选定前$k$个线性无关列后，下一列有$q^n-q^k$种选择。因此
\[
|M_n(F)|=q^{n^2},\qquad
|\GL_n(F)|=\prod_{k=0}^{n-1}(q^n-q^k).
\]
\textbf{出处：}现有笔记抽象代数第一章习题44的逐列计数（\path{chapter1-exercises.tex:454}）。原题取$q=p$为素数，本条将同一证明适用的任意有限域条件完整写明；列满秩等价于可逆是此解所用的线性代数工具。
\end{reviewprop}

\begin{reviewprop}[连续函数环所用的分析工具（补充）]\label{tool:g-analysis}
\textbf{紧致性：}闭区间$[0,1]$的任意相对开覆盖都有有限子覆盖。
\textbf{连续性：}有限个实值连续函数的和、积以及平方仍连续；若连续函数$g$在$[0,1]$上处处非零，则$1/g$连续。
\textbf{非零邻域：}若$g(t)\ne0$，则存在$\delta>0$，使$g(x)\ne0$对所有$x\in(t-\delta,t+\delta)\cap[0,1]$成立。
\textbf{局部非零函数：}任意非空相对开区间$J\subseteq[0,1]$内，可取$c\in(0,1)$、$\varepsilon>0$使$[c-\varepsilon,c+\varepsilon]\subseteq J$。连续函数
\[
g(x)=\max\{0,1-|x-c|/\varepsilon\}
\]
满足$g(c)=1$且在该小闭区间外为零。
\textbf{出处：}分析工具补充，系现有笔记抽象代数第一章习题43证明实际使用的紧致性、连续倒数、非零邻域与三角形函数（\path{chapter1-exercises.tex:440}）。不标作代数正文原定理，也不声称这些分析定理在代数笔记中已有完整证明。
\end{reviewprop}

\begin{reviewprop}[正多边形的几何二面体群]\label{tool:g-dihedral}
设$n\geq3$。正$n$边形的全部平面对称变换在复合下成群，称$D_n$。令$T$为绕中心旋转$2\pi/n$，$S$为关于一条对称轴的反射，则
\[
D_n=\{T^k,ST^k:0\leq k<n\},\quad |D_n|=2n,\quad
\operatorname{ord}(T)=n,\quad S^2=e,\quad STS=T^{-1}.
\]
该群不交换：否则$ST=TS$导致$T^2=e$，与$n\geq3$矛盾。此处$D_n$采用阶为$2n$的记号约定。
\textbf{出处：}现有笔记抽象代数第一章第3节“群的例子”之“图形的对称群与二面体群”，对称群命题、二面体群定义及其后的阶数命题（\path{chapter1.tex:154,164,170}）。
\end{reviewprop}

\begin{reviewtheorem}[由方程可解性判定群]\label{tool:g-group-construction}
设非空集合$G$上有结合的二元运算。如果对每一对$a,b\in G$，方程$ax=b$和$ya=b$在$G$中均有解，则$G$是群。条件不要求$G$有限，也不要求预先假设解唯一；单位元与逆元须由上述可解性构造，不能在构造前使用群的消去律。
\textbf{出处：}现有笔记抽象代数第一章习题4“可解方程的结合系统”（\path{chapter1-exercises.tex:34}）；本条是该习题结论的编号陈述，其构造性证明保留在详解中。
\end{reviewtheorem}

% 工具编号由主文件的 reviewtheorem/reviewprop/reviewlemma 共享计数器统一管理。
\subsection{群论的作用、生成与可解性工具}
以下条目完整列出第二章详解实际使用的结论。所标行号指原笔记源文件，不是教材定理编号；原习题解答中的补充结论也明确标出，不冒称为教材已证明的定理。各条采用本复习册统一的工具编号，题目中的引用可直接跳转至相应条目。

\begin{reviewtheorem}[对换与3-轮换的生成结论]\label{tool:a-permutation-generation}
设$n\ge2$。$S_n$中的每个置换都是有限个对换的乘积，并且
\[
S_n=\langle(12),(13),\ldots,(1n)\rangle
=\langle(12),(23),\ldots,(n-1\ n)\rangle.
\]
设$n\ge3$。$A_n$中的每个偶置换都是有限个3-轮换的乘积，因此$A_n$由全部3-轮换生成。

\textbf{原文出处：}原笔记第二章“单群与交错群的单性”中的引理，\path{chapters/Part 2/chapter2.tex} 第362--381行；两组指定对换的生成结论是原习题10及其补充证明，\path{chapters/Part 2/chapter2-exercises.tex} 第201--219行。证明展开见本册第二章习题10。
\end{reviewtheorem}

\begin{reviewprop}[轮换分解、置换的阶、共轭公式与奇偶性]\label{tool:a-cycle-conjugation}
每个$\sigma\in S_n$都可写成支集两两不相交的轮换之积；在忽略轮换排列次序、每个轮换的循环起点以及省略的不动点后，该分解唯一。
若非平凡轮换长度为$d_1,\ldots,d_s$，则
\[
\operatorname{ord}(\sigma)=\lcm(d_1,\ldots,d_s).
\]
恒等置换的阶为1。这个阶公式是分解定理的补充推论：不相交轮换分别作用于各自支集，所以对任意正整数$m$，$\sigma^m=e$当且仅当每个$d_i$都整除$m$，最小这样的正整数就是最小公倍数。

设$\pi\in S_n$，$a_1,\ldots,a_r$为互不相同的点，$2\le r\le n$。则
\[
\pi(a_1\cdots a_r)\pi^{-1}
=(\pi(a_1)\cdots\pi(a_r)).
\]
$r$-轮换的符号为$(-1)^{r-1}$。符号映射$\sgn:S_n\to\{1,-1\}$为群同态，其核为$A_n$；支集不相交的置换互相交换。特别地，奇数长度的轮换是偶置换。

\textbf{原文出处：}不相交轮换分解定理及其证明见原笔记第一章第3节“群的例子”，\path{chapters/Part 2/chapter1.tex} 第222--234行；置换的阶为轮换长度最小公倍数是本条明确补出的推论，不标作原笔记独立阶公式定理。共轭公式见第二章“单群与交错群的单性”的轮换共轭命题，\path{chapters/Part 2/chapter2.tex} 第383--410行；轮换的对换分解见同文件第366--370行，奇偶性和不相交支集的使用见原习题12、40，\path{chapters/Part 2/chapter2-exercises.tex} 第258--276、844--857行。这里将上述已使用的置换规则合并列出。
\end{reviewprop}

\begin{reviewlemma}[逐次附加点的3-轮换引理（补充）]\label{tool:a-triple-attach}
设$s\ge1$，$\tau_1,\ldots,\tau_s$是$S_n$中的3-轮换，$n\ge3$。记$U_j$为前$j$个轮换所涉及的全部点。若对每个$j\ge2$，$\tau_j$的三个点中至少有一点属于$U_{j-1}$，且$U_s=\{1,\ldots,n\}$，则
\[
\langle\tau_1,\ldots,\tau_s\rangle=A_n.
\]
其中使用的单步结论是：若一个子群已包含点集$U$（$|U|\ge3$）上的全部偶置换，并含$(abz)$，$a,b\in U$、$a\ne b$、$z\notin U$，则它包含$U\cup\{z\}$上的全部偶置换。这里“点集上的偶置换”均延拓为在其余点上恒等的置换。

\textbf{原文出处：}原习题12之前的补充引理及证明，\path{chapters/Part 2/chapter2-exercises.tex} 第233--252行；不是教材习题页给出的定理。完整证明见本册第二章习题12，保留原笔记先处理一个新点、再由四点交错群处理两个新点的顺序。
\end{reviewlemma}

\begin{reviewprop}[商子群的代表元判据]\label{tool:a-quotient-membership}
设$N\trianglelefteq G$、$H\le G$且$N\le H$。则$N\trianglelefteq H$，$H/N$是$G/N$的子群，并且对每个$g\in G$都有
\[
gN\in H/N\iff g\in H.
\]
因此若$N\le H_1\cap H_2$，则在$G/N$中有实际相等关系
\[
(H_1/N)\cap(H_2/N)=(H_1\cap H_2)/N.
\]

\textbf{原文出处：}原笔记第二章“群的同态定理”的子群对应定理，\path{chapters/Part 2/chapter2.tex} 第10--63行；代表元判据及交的应用见原习题15，\path{chapters/Part 2/chapter2-exercises.tex} 第331--346行。证明见本册第二章习题15。
\end{reviewprop}

\begin{reviewtheorem}[轨道、稳定子与共轭子群数]\label{tool:a-orbit-stabilizer}
设群$G$作用于集合$X$，即$e\cdot x=x$、$(gh)\cdot x=g\cdot(h\cdot x)$。对$x\in X$定义轨道$G\cdot x=\{g\cdot x:g\in G\}$与稳定子$G_x=\{g\in G:g\cdot x=x\}$。$G_x$为子群，各轨道将$X$分为不交并，且
\[
\{gG_x:g\in G\}\longrightarrow G\cdot x,
\qquad gG_x\longmapsto g\cdot x
\]
为双射。若$y=g\cdot x$，则$G_y=gG_xg^{-1}$；若$G$有限，则$|G\cdot x|=[G:G_x]$。

特别地，对$H\le G$定义$N_G(H)=\{g:gHg^{-1}=H\}$，它为子群且包含$H$。当$G$有限时，不同共轭子群$gHg^{-1}$的个数为
\[
[G:N_G(H)]\le[G:H].
\]

\textbf{原文出处：}原笔记第二章“群在集合上的作用”的轨道--稳定子定理，\path{chapters/Part 2/chapter2.tex} 第886--930行；正规化子的定义及验证见第976--990行；共轭子群数的应用见原习题30，第611--627行。
\end{reviewtheorem}

\begin{reviewtheorem}[类方程与有限$p$群的中心]\label{tool:a-class-equation}
设$G$为有限群，$Z(G)=\{z\in G:zg=gz\ (g\in G)\}$，并定义中心化子$C_G(x)=\{g\in G:gx=xg\}$。若$x_1,\ldots,x_t$分别代表全部非中心共轭类，则
\[
|G|=|Z(G)|+\sum_{i=1}^t[G:C_G(x_i)].
\]
若$p$为素数、$|G|=p^a$且$a\ge1$，则$p\mid|Z(G)|$，从而$Z(G)\ne\{e\}$。中心$Z(G)$总是$G$的正规交换子群。

\textbf{原文出处：}原笔记第二章“群在集合上的作用”的类方程定理及证明，\path{chapters/Part 2/chapter2.tex} 第953--974行；中心的定义在第757--758行，中心作为正规交换子群用于原习题44，第904--916行。
\end{reviewtheorem}

\begin{reviewtheorem}[西罗第一定理：素数幂阶子群的存在]\label{tool:a-sylow-existence}
设$G$为有限群、$p$为素数，$k\ge0$为整数，且$p^k\mid|G|$。则$G$存在一个阶为$p^k$的子群。特别地，若$|G|=p^am$、$p\nmid m$，则存在$p^a$阶子群，称为西罗$p$子群；阶为$p$的幂的群称为$p$群。

\textbf{原文出处：}原笔记第二章“西罗定理”的第一定理及定义，\path{chapters/Part 2/chapter2.tex} 第1020--1046行。$k=0$时取平凡子群。
\end{reviewtheorem}

\begin{reviewtheorem}[西罗第二定理：包含与共轭]\label{tool:a-sylow-conjugacy}
设$G$为有限群、$p$为素数，$P$是$G$的一个西罗$p$子群。对任意$p$子群$H\le G$，存在$g\in G$使
\[
H\le gPg^{-1}.
\]
因此每个$p$子群都包含于一个西罗$p$子群，且任意两个西罗$p$子群互相共轭。

\textbf{原文出处：}原笔记第二章“西罗定理”的第二定理及证明，\path{chapters/Part 2/chapter2.tex} 第1049--1066行。
\end{reviewtheorem}

\begin{reviewtheorem}[西罗第三定理：个数与正规性]\label{tool:a-sylow-count}
设$G$为有限群、$p$为素数，$|G|=p^am$、$p\nmid m$，$P$为一个西罗$p$子群，$n_p$为全部西罗$p$子群的个数。则
\[
n_p=[G:N_G(P)],\qquad n_p\mid m,\qquad n_p\equiv1\pmod p.
\]
并且$n_p=1$当且仅当$P\trianglelefteq G$。这些结论也允许$a=0$，此时唯一的西罗$p$子群为平凡群。

\textbf{原文出处：}原笔记第二章“西罗定理”的第三定理及证明，\path{chapters/Part 2/chapter2.tex} 第1069--1096行。
\end{reviewtheorem}

\begin{reviewprop}[素数平方阶群交换（原习题结论）]\label{tool:a-p-square-abelian}
设$p$为素数、$G$为有限群，$|G|=p^2$。则$G$交换，而且恰有两种可能的同构类型：循环群$C_{p^2}$与$C_p\times C_p$。其中$C_m$表示$m$阶循环群。

\textbf{原文出处：}原习题20的题目及补充证明，\path{chapters/Part 2/chapter2-exercises.tex} 第422--441行。类方程与中心商群的交换性论证在本册第二章习题43第二步展开；本册使用此条时只需其交换性部分。
\end{reviewprop}

\begin{reviewprop}[$p^2q$阶群存在正规西罗子群（补充）]\label{tool:a-p-square-q-normal}
设$p,q$为不同素数、$G$为有限群，$|G|=p^2q$。则$G$的西罗$p$子群或西罗$q$子群至少有一种是唯一的，因而$G$至少有一个非平凡真正规西罗子群。

\textbf{原文出处：}原习题25的补充结论与证明，\path{chapters/Part 2/chapter2-exercises.tex} 第510--528行；不是正文另编号的定理。证明见本册第二章习题43第一步，包括$p=2,q=3$的12阶计数情况。
\end{reviewprop}

\begin{reviewprop}[换位子群、同态像与导出列]\label{tool:a-derived-functor}
对任意群$G$，令$[x,y]=x^{-1}y^{-1}xy$，$G^{(1)}=\langle[x,y]:x,y\in G\rangle$，并递归定义$G^{(0)}=G$、$G^{(i+1)}=(G^{(i)})^{(1)}$。则$G^{(1)}\trianglelefteq G$，且$G$交换当且仅当$G^{(1)}=\{e\}$。

若$H\le G$，则对每个$i\ge0$有$H^{(i)}\le G^{(i)}$。若$\varphi:G\to Q$为群同态，则
\[
\varphi([x,y])=[\varphi(x),\varphi(y)],\qquad
\varphi(G^{(i)})=(\varphi(G))^{(i)}.
\]
特别地，当$\varphi$满射时，$\varphi(G^{(i)})=Q^{(i)}$。因此对$N\trianglelefteq G$，
\[
G/N\text{交换}\iff G^{(1)}\le N.
\]

\textbf{原文出处：}原笔记第二章“群的同态定理”的换位子群定义与交换商群判据，\path{chapters/Part 2/chapter2.tex} 第161--193行；“可解群”的导出列定义与子群、同态像证明分别在第523--539、695--719行。原习题42的补充证明在\path{chapters/Part 2/chapter2-exercises.tex} 第874--889行。
\end{reviewprop}

\begin{reviewprop}[可解群的子群与同态像]\label{tool:a-solvable-closure}
称群$G$可解，是指存在整数$r\ge0$使$G^{(r)}=\{e\}$，导出列按\toolref{a-derived-functor}定义。若$G$可解，则任意子群$H\le G$和任意同态像$\varphi(G)$均可解。特别地，每个商群$G/N$（$N\trianglelefteq G$）可解。可解性在同构下保持；交换群均可解。

\textbf{原文出处：}原笔记第二章“可解群”的定义及说明，\path{chapters/Part 2/chapter2.tex} 第523--539行，以及关于子群、同态像和扩张的命题及证明，第695--719行；本册第二章习题42保留原证明顺序。
\end{reviewprop}

\begin{reviewprop}[可解群的扩张性质]\label{tool:a-solvable-extension}
设$G$为群，$N\trianglelefteq G$。若$N$与$G/N$均可解，则$G$可解。更具体地，若$N^{(s)}=\{e\}$、$(G/N)^{(r)}=\{N\}$，其中$r,s\ge0$，则
\[
G^{(r)}\le N,\qquad G^{(r+s)}=\{e\}.
\]
此结论不要求$G$有限。特别地，若$N$和$G/N$均交换，则$G^{(2)}=\{e\}$。

\textbf{原文出处：}原笔记第二章“可解群”中关于子群、同态像和扩张的命题，\path{chapters/Part 2/chapter2.tex} 第695--724行；扩张部分证明见第716--724行。本册第二章习题43第三步按原方法展开，习题44再次应用。
\end{reviewprop}

\begin{reviewprop}[可解群的有限直积]\label{tool:a-solvable-product}
设$A,B$为群，其直积按逐分量乘法定义。对所有$i\ge0$有
\[
(A\times B)^{(i)}=A^{(i)}\times B^{(i)}.
\]
因此若$A^{(r)}=\{e_A\}$、$B^{(s)}=\{e_B\}$，则$A\times B$的第$\max\{r,s\}$个导群平凡，$A\times B$可解。反复使用该结论，任意有限多个可解群的直积可解，其任意子群也可解。

\textbf{原文出处：}原习题45补充证明中的直积导群计算，\path{chapters/Part 2/chapter2-exercises.tex} 第919--935行；直积的逐分量运算定义见原笔记第二章“群的直和”，\path{chapters/Part 2/chapter2.tex} 第1124--1133行。完整计算见本册第二章习题45。
\end{reviewprop}

\begin{reviewtheorem}[内自同构群与中心商群]\label{tool:a-inner-automorphisms}
设$G$为群。对$g\in G$，令$c_g:G\to G$、$x\mapsto gxg^{-1}$，则$c_g$为自同构。全体内自同构组成群$\operatorname{Inn}(G)$。映射
\[
G\twoheadrightarrow\operatorname{Inn}(G),\qquad g\longmapsto c_g
\]
为满同态，其核为$Z(G)$，故$G/Z(G)\cong\operatorname{Inn}(G)$。此外，$\operatorname{Inn}(G)\trianglelefteq\Aut(G)$，且对每个$\alpha\in\Aut(G)$有
\[
\alpha c_g\alpha^{-1}=c_{\alpha(g)}.
\]
若$G$有限且$Z(G)=\{e\}$，则恰有$|G|$个内自同构。中心平凡且全部自同构均为内自同构的群称为完全群。

\textbf{原文出处：}原笔记第二章“群的自同构群”的内自同构定理，\path{chapters/Part 2/chapter2.tex} 第757--791行；完全群的定义在第809--812行。
\end{reviewtheorem}

\begin{reviewprop}[自同构保持的元素与子群性质]\label{tool:a-automorphism-invariants}
设$\alpha\in\Aut(G)$。则$\alpha(e)=e$，且每个元素$g$与$\alpha(g)$的阶相同（包括无限阶）。若$H\le G$，则$\alpha(H)\le G$，$\alpha$限制为$H$到$\alpha(H)$的同构；若$G$有限，则它把西罗$p$子群互相置换。并且
\[
\alpha(N_G(H))=N_G(\alpha(H)).
\]
若一个自同构固定某个生成集的每个元素，则它为恒等自同构。对由元素阶、子群阶等这些不变性质确定的有限集合或子群族，自同构给出其上的置换；这个置换作用是否单射仍须另证。

\textbf{原文出处：}原习题49用元素阶、西罗子群及其正规化子构造自同构的作用，\path{chapters/Part 2/chapter2-exercises.tex} 第991--1018行；上述性质是该解答所用的同构定义直接推论，这里作为补充工具完整列出，不冒称为正文独立编号命题。定义和内自同构的共轭计算见原笔记第二章“群的自同构群”，\path{chapters/Part 2/chapter2.tex} 第729--731、784--791行。
\end{reviewprop}

\begin{reviewprop}[循环群的自同构]\label{tool:a-cycle-automorphisms}
设$C_m=\langle a\rangle$为$m$阶循环群，$m\ge2$。每个自同构由$a$的像唯一决定，且$a$必须映为一个生成元。因此
\[
\Aut(C_m)=\{\alpha_k:\alpha_k(a^t)=a^{kt},\ \gcd(k,m)=1\},
\qquad
\Aut(C_m)\cong(\Z/m\Z)^\times.
\]
这里右边是模$m$的可逆剩余类乘法群。特别地，对素数$p$，有$|\Aut(C_p)|=p-1$。无限循环群的自同构只有恒等与取逆两种，平凡群只有恒等自同构。

\textbf{原文出处：}原笔记第二章“群的自同构群”关于$\Aut(\Z)$与$\Aut(C_m)$的命题及证明，\path{chapters/Part 2/chapter2.tex} 第733--754行。
\end{reviewprop}

\begin{reviewlemma}[指数2交换群与基向量交换（补充线性代数工具）]\label{tool:a-binary-basis-swap}
设$G$为有限交换群，且$x^2=e$对每个$x\in G$成立。把群运算改记为加法、$e$记为0，规定$0x=0$、$1x=x$，则$G$成为$\F_2$上的有限维向量空间。它存在一组基$v_1,\ldots,v_d$，每个元素唯一表示为$\sum_i\varepsilon_iv_i$，$\varepsilon_i\in\F_2$，因此$|G|=2^d$。

若$|G|>2$，则$d\ge2$。把$v_1,v_2$交换而固定其余基向量，可唯一线性延拓为一个可逆线性映射；该映射也是$G$的一个非恒等群自同构。

\textbf{原文出处：}原习题47第三种情况的补充证明，\path{chapters/Part 2/chapter2-exercises.tex} 第964--969行。向量空间结构、有限基与线性延拓的具体使用在本册第二章习题47中展开；这里是该题原方法所用的补充工具，并非本章教材另列的定理。
\end{reviewlemma}

\begin{reviewprop}[一般线性群中的单位上三角西罗子群（补充）]\label{tool:a-gl-sylow}
设$p$为素数、$n\ge1$，$G=\GL_n(\F_p)$。令$U$为全体对角线全1的上三角矩阵，$B$为全体可逆上三角矩阵，则$U,B$均为$G$的子群，且
\[
|U|=p^{n(n-1)/2},\qquad
|B|=(p-1)^np^{n(n-1)/2},\qquad
N_G(U)=B.
\]
结合\toolref{g-basis-count}的可逆矩阵计数公式
\[
|G|=p^{n(n-1)/2}\prod_{i=1}^n(p^i-1),
\]
可知$U$是西罗$p$子群，所有西罗$p$子群均为$U$的共轭，个数为
\[
n_p=[G:B]=\frac{\prod_{i=1}^n(p^i-1)}{(p-1)^n}.
\]
$n=1$时$U=\{1\}$、$B=G$，上述结论仍成立。

\textbf{原文出处：}原习题36及其补充解答，\path{chapters/Part 2/chapter2-exercises.tex} 第733--786行。正规化子的证明使用公共不动空间递归恢复标准子空间链，本册第二章习题36保留并展开该原方法；基础的$\GL_n$阶数见\toolref{g-basis-count}的来源。
\end{reviewprop}

\begin{reviewprop}[中心商群循环则原群交换]\label{tool:a-center-quotient}
设$G$为任意群，不要求有限。若商群$G/Z(G)$循环，则$G$交换。具体地，若$G/Z(G)=\langle aZ(G)\rangle$，则每个元素都可写成$a^mz$，其中$m\in\Z$、$z\in Z(G)$，任意两个这样的元素互相交换。因此中心商群一旦循环，就只能为平凡群。

\textbf{原文出处：}原习题39的补充证明，\path{chapters/Part 2/chapter2-exercises.tex} 第823--838行；在原习题20、37、46中也使用此结论。完整证明见本册第二章习题39。
\end{reviewprop}

\begin{reviewtheorem}[凯莱定理与有限群的正则嵌入]\label{tool:a-cayley}
每个群$G$都同构于集合$G$上的一个置换群。具体地，对$g\in G$定义左乘置换$L_g(x)=gx$，则
\[
G\longrightarrow\operatorname{Sym}(G),\qquad g\longmapsto L_g
\]
为单射群同态。特别地，若$|G|=m<\infty$，给$G$的元素编号后得到$G\hookrightarrow S_m$。

\textbf{原文出处：}原笔记第一章“群的同构”中的凯莱定理及左正则表示证明，\path{chapters/Part 2/chapter1.tex} 第406--469行；原笔记第二章“群在集合上的作用”的左平移作用见\path{chapters/Part 2/chapter2.tex} 第857行；本册第二章习题40展开该嵌入。
\end{reviewtheorem}

% !TEX root = main.tex
\subsection{非交换环、逆元与商环工具}

\begin{reviewprop}[环的零、负元与分配律]\label{tool:r-ring-laws}
环的加法是交换群，故可在加法等式两侧消去同一加数。乘法结合并对加法左右分配，不预设乘法交换。对任意$a,b$，
\[
a0=0a=0,\quad a(-b)=(-a)b=-(ab),\quad(-a)(-b)=ab.
\]
分配展开和提取公因子时，须保留乘法次序；一般不能把$ab$换成$ba$。

\textbf{出处：}现有代数笔记抽象代数第一章第8节“环、子环”，环定义与零、负元运算命题；源文件\path{chapter1.tex}第779--827行。
\end{reviewprop}

\begin{reviewprop}[双侧逆元唯一与左右逆相等]\label{tool:r-unit-laws}
在幺环$R$中，若$uv=vu=1$，则$u$为单位、$v=u^{-1}$，且双侧逆元唯一。若$b u=1$、$u c=1$，则$b=b(uc)=(bu)c=c$，所以一个左逆和一个右逆同时存在时必相等。仅给定$uv=1$不能在无限环中直接断言$vu=1$。

给定$uv=1$，若$xu=0$，则$x=(xu)v=0$；并且对整数$0\le i\le j$，
\[
u^jv^j=1,\qquad u^iv^j=v^{j-i}.
\]
这些幂等式由结合律逐对消去中间的$uv$，不需要$uv=vu$。

\textbf{出处：}现有笔记第一章第9节“各种特殊类型的环”，单位与逆元唯一性命题，\path{chapter1.tex}第952--972行；幂消去是第三章习题3原解中使用并在本册题内归纳验证的步骤。
\end{reviewprop}

\begin{reviewprop}[环同态保持的等式与单位]\label{tool:r-homomorphism}
环同态$\eta:R\to S$保加法和乘法，因此对整数$k$及正整数$m$有
\[
\eta(0)=0,\quad\eta(-a)=-\eta(a),\quad
\eta(ka)=k\eta(a),\quad\eta(a^m)=\eta(a)^m.
\]
因而幂零元、幂等元的像仍为相应元素。若$R,S$均为幺环且$\eta(1_R)=1_S$，则单位$u$的像是单位，且$\eta(u)^{-1}=\eta(u^{-1})$，不要求$\eta$满射。一般环同态不预设保幺；“单位的像是单位”必须保留保幺条件。零因子要求元素本身及消去它的另一个因子都非零，这两个非零条件不由一般同态保持。

\textbf{出处：}第一章第10节“环的同态、理想”，同态定义及基本性质，\path{chapter1.tex}第1150--1198行；正幂由保乘法归纳。关于幂零元、幂等元和单位的结论在原第三章习题8及本册该题中实际证明，列在此处便于定位，不据此省略原题证明。
\end{reviewprop}

\begin{reviewprop}[整数商环与自然满同态]\label{tool:r-quotients}
对正整数$m$，$\Z/m\Z$以剩余类逐项加、乘构成交换幺环，单位元是$\overline1_m$，映射$k\mapsto\overline k_m$为保幺满同态。若正整数$d\mid m$，则
\[
\Z/m\Z\longrightarrow\Z/d\Z,\qquad\overline k_m\longmapsto\overline k_d
\]
良定义且为保幺满同态。$m\ge2$时，$\overline a_m$为单位当且仅当$\gcd(a,m)=1$；$\Z/m\Z$为整环当且仅当$m$为素数，此时也是域。整数环$\Z$本身是整环，单位只有$1,-1$。

\textbf{出处：}第一章第11节“商环”，商环与自然同态以及整数模$n$环性质，\path{chapter1.tex}第1250--1388行；整数环及其零因子、单位见第一章第8、9节。$d\mid m$保证映射良定义这一检查仍在第三章习题8各反例中保留。
\end{reviewprop}

\setcounter{secnumdepth}{2}
\clearpage
% !TEX root = main.tex
\section{第零章：集合与整数（补充证明）}
本章题面据教材书页20--22核对。二元关系是集合$S\times S$的子集，记$x\sim y$表示$(x,y)$属于该关系。自反性指每个$x\in S$满足$x\sim x$；对称性指$x\sim y$推出$y\sim x$；传递性指$x\sim y$且$y\sim z$推出$x\sim z$。三者同时成立时称为等价关系。

\subsection{习题4：判断关系的自反性、对称性与传递性}\label{exercise:0:4}
\paragraph{题目}判断以下关系是否具有自反性、对称性、传递性：
\begin{enumerate}[label=(\roman*)]
\item $S=\R$，$x-y$为$2\pi$的整数倍；
\item $S=\R$，$x\geq y$；
\item $S=\R$，$x>y$；
\item $S$仅含一个元素，关系为$x>y$。
\end{enumerate}
\paragraph{考点与思路}证明某性质须对任意元素成立；否定只需一个反例。空关系的传递性与对称性是条件命题，前提不能成立时仍成立。\textbf{原文依据：}第零章第1节等价关系的三条定义，书页2；本题无附带证明，按原定义逐项补证。
\paragraph{所用结论}\toolref{n-relations}。
\begin{proof}[详细证明]
(i) 因$x-x=2\pi\cdot0$，自反性成立。若$x-y=2\pi k$，$k\in\Z$，则$y-x=2\pi(-k)$，故对称性成立。若另有$y-z=2\pi l$，则$x-z=2\pi(k+l)$，故传递性成立。

(ii) 任意$x$满足$x\geq x$，所以自反。取$x=1,y=0$，有$x\geq y$却没有$y\geq x$，所以不对称。由$x\geq y$及$y\geq z$得到$x\geq z$，所以传递。

(iii) $x>x$对每个实数都不成立，故不自反。取$1>0$，反向$0>1$不成立，故不对称。若$x>y>z$，则$x>z$，故传递。

(iv) 设$S=\{a\}$。因$a>a$不成立，关系是空关系，且不自反。不存在满足$x>y$的$x,y\in S$，因此不存在违反对称性或传递性的实例，这两个性质均成立。
\end{proof}
\paragraph{复习提示}依次为“全成立”“自反、传递”“仅传递”“对称、传递”。最后一项不能因关系为空而把所有性质一律判为不成立。

\subsection{习题5：差与和所定义的关系}\label{exercise:0:5}
\paragraph{题目}在$\R$上分别以(i) $|x-y|=1$；(ii) $|x-y|\leq1$；(iii) $x-y=1$；(iv) $x+y=1$定义关系，判断其自反性、对称性和传递性。
\paragraph{考点与思路}把三个性质逐项代入，注意距离的三角不等式只能给出上界$2$，不能自动保持上界$1$。\textbf{原文依据：}第零章第1节等价关系的三条定义，书页2；本题无附带证明，按原定义逐项补证。
\paragraph{所用结论}\toolref{n-relations}。
\begin{proof}[详细证明]
(i) $|x-x|=0\ne1$，不自反。由$|x-y|=|y-x|$，对称。取$x=0,y=1,z=0$，前两个距离均为$1$，而$|x-z|=0$，不传递。

(ii) $|x-x|=0\leq1$，自反；距离对称，故关系对称。取$x=0,y=1,z=2$，有$|x-y|=|y-z|=1$，而$|x-z|=2>1$，不传递。

(iii) $x-x=0$，不自反。取$x=1,y=0$，则$x-y=1$而$y-x=-1$，不对称。取$x=2,y=1,z=0$，则$x-y=y-z=1$而$x-z=2$，不传递。

(iv) 要自反须所有实数均满足$2x=1$，取$x=0$即否定。交换加数不改变和，因此对称。取$x=0,y=1,z=0$，有$x+y=y+z=1$而$x+z=0$，不传递。
\end{proof}
\paragraph{复习提示}依次为“仅对称”“自反、对称”“全不成立”“仅对称”。反例必须同时满足传递性中的两个前提。

\subsection{习题6：自然数加法、乘法的结合律与交换律}\label{exercise:0:6}
\paragraph{题目}证明自然数的加法和乘法分别满足结合律和交换律。
\paragraph{考点与思路}\textbf{原文依据：}第零章第2节的递归定理、加法与乘法递归定义，书页5--6。原文将运算律证明省略，以下从同一递归定义出发，用数学归纳法补全。与原文一致取$\N=\{0,1,2,\ldots\}$，写$S(n)=n+1$表示原文的后继$n^+$，使用递归定义
\[
a+0=a,\quad a+S(b)=S(a+b),\quad a\cdot0=0,\quad aS(b)=ab+a.
\]
\paragraph{所用结论}\toolref{n-recursion}、\toolref{n-induction}。
\begin{proof}[详细证明]
\textbf{第一步：加法的两个准备等式。}对$b$归纳，$0+0=0$，且若$0+b=b$，则$0+S(b)=S(0+b)=S(b)$。所以$0+b=b$。

同样对$b$归纳证明$S(a)+b=S(a+b)$：$b=0$时成立；若对$b$成立，则
\[
S(a)+S(b)=S(S(a)+b)=S(S(a+b))=S(a+S(b)).
\]
\textbf{第二步：加法交换律。}固定$a$，对$b$归纳。$a+0=0+a$。若$a+b=b+a$，则
\[
a+S(b)=S(a+b)=S(b+a)=S(b)+a.
\]
因此$a+b=b+a$。

\textbf{第三步：加法结合律。}固定$a,b$，对$c$归纳。$c=0$时$(a+b)+0=a+(b+0)$。若对$c$成立，则
\begin{align*}
(a+b)+S(c)&=S((a+b)+c)\\
&=S(a+(b+c))=a+(b+S(c)).
\end{align*}
故$(a+b)+c=a+(b+c)$。

\textbf{第四步：先证左分配律。}固定$a,b$，对$c$归纳。$a(b+0)=ab=ab+a0$。若$a(b+c)=ab+ac$，则
\begin{align*}
a(b+S(c))&=aS(b+c)=a(b+c)+a\\
&=(ab+ac)+a=ab+(ac+a)=ab+aS(c).
\end{align*}
这一步只用了递归定义和已证加法结合律，并未使用乘法交换律。

\textbf{第五步：乘法交换律。}先对$b$归纳得$0b=0$。再对$b$归纳证明$S(a)b=ab+b$。初始$b=0$成立；归纳步中
\begin{align*}
S(a)S(b)&=S(a)b+S(a)=(ab+b)+(a+1)\\
&=(ab+a)+(b+1)=aS(b)+S(b),
\end{align*}
其中加数的换序只使用已经证明的加法交换律与结合律。

现对$b$归纳：$a0=0=0a$；若$ab=ba$，则
\[
aS(b)=ab+a=ba+a=S(b)a.
\]
故$ab=ba$。

\textbf{第六步：乘法结合律。}固定$a,b$，对$c$归纳。$c=0$时两边为$0$；若$(ab)c=a(bc)$，则用第四步得
\begin{align*}
(ab)S(c)&=(ab)c+ab=a(bc)+ab\\
&=a(bc+b)=a(bS(c)).
\end{align*}
所以$(ab)c=a(bc)$。四项结论全部得证。
\end{proof}
\paragraph{复习提示}不要在证明乘法交换律之前，未经证明就交换乘数。这里先证左分配律作为引理，避免各运算律之间循环引用。

\subsection{习题7：乘法对加法的分配律}\label{exercise:0:7}
\paragraph{题目}证明自然数乘法对加法的分配律成立。
\paragraph{考点与思路}以加法的第二个变量为归纳变量；右分配律再由乘法交换律推出。\textbf{原文依据：}第零章第2节的乘法递归定义与原文省略证明的分配律，书页5--6；以下为在该定义下的补充证明。
\paragraph{所用结论}\toolref{n-recursion}、\toolref{n-induction}。另使用前题已证明的加法结合律与乘法交换律（第零章习题6）。
\begin{proof}[详细证明]
固定$a,b\in\N$，对$c$证明$a(b+c)=ab+ac$。$c=0$时两边均为$ab$。设对$c$已成立，则
\begin{align*}
a(b+(c+1))&=a((b+c)+1)\\
&=a(b+c)+a\\
&=(ab+ac)+a\\
&=ab+(ac+a)=ab+a(c+1).
\end{align*}
第一、第四步用加法结合律，第二、末一步用乘法递归定义。故左分配律成立。

由习题6的乘法交换律，
\[
(a+b)c=c(a+b)=ca+cb=ac+bc.
\]
因而右分配律也成立。
\end{proof}
\paragraph{复习提示}归纳步须写清楚如何把$c+1$化归$c$，不能仅说“由归纳法成立”。

\subsection{习题8：自然数乘法的消去律}\label{exercise:0:8}
\paragraph{题目}若$ax=ay$且$a\ne0$，证明$x=y$。
\paragraph{考点与思路}使用自然数的次序：若$x<y$，则$y=x+t$，其中$t$为正整数，乘以正整数后差仍为正。\textbf{原文依据：}第零章第2节自然数次序以$a+t=b$定义，书页6，以及习题7所补的分配律；原书未附本题证明，以下据这些原有工具补证。
\paragraph{所用结论}\toolref{n-order}、\toolref{n-recursion}、\toolref{n-induction}，并用本章习题7已证的分配律。
\begin{proof}[详细证明]
自然数乘法递归定义表明：对正整数$a$，$a\cdot1=a>0$；若$at>0$，则$a(t+1)=at+a>0$。所以对所有正整数$t$都有$at>0$。

若$x\ne y$，自然数的三歧性保证$x<y$或$y<x$。若$x<y$，取正整数$t$使$y=x+t$。由分配律，
\[
ay=a(x+t)=ax+at>ax,
\]
与$ax=ay$矛盾。若$y<x$，同理得$ax>ay$，也矛盾。因此$x=y$。
\end{proof}
\paragraph{复习提示}$a\ne0$不可删去，$0\cdot1=0\cdot2$即为反例。此证明没有先假设分式或除法已经定义。

\subsection{习题13：两类大整数的最大公因数与素数无穷性}\label{exercise:0:13}
\paragraph{题目}设$a>1$为正整数，$m>n>0$，$m,n\in\Z$。证明
\[
\gcd(a^{2^m}+1,a^{2^n}+1)=1\text{或}2,
\]
并由此证明有无穷多个素数。
\paragraph{考点与思路}模共同因子把$a^{2^n}$替换成$-1$。为推出素数无穷性，选择偶数$a$，使最大公因数必为$1$。\textbf{原文依据：}第零章第3节整除与素因子、算术基本定理（书页10--11）和第4节同余运算法则（书页12）；原书本题没有提供证明或提示，以下为使用本章工具的补充证明，不标作原书原证。
\paragraph{所用结论}\toolref{n-gcd}、\toolref{n-prime-factor}、\toolref{n-congruence}。
\begin{proof}[详细证明]
设$d$为上述两数的最大公因数，则$a^{2^n}\equiv-1\pmod d$。因$m>n$，$2^{m-n}$为偶数，故
\[
a^{2^m}=(a^{2^n})^{2^{m-n}}\equiv(-1)^{2^{m-n}}=1\pmod d.
\]
另一方面$d\mid a^{2^m}+1$，所以同一数也同余于$-1$。相减得$d\mid2$，即$d=1$或$2$。

更具体地，若$a$为偶数，两数均为奇数，故$d=1$；若$a$为奇数，两数均为偶数，故$d=2$。

令$a=2$。对每个$n\geq1$，整数$F_n=2^{2^n}+1>1$，故至少有一个素因子$p_n$。当$m\ne n$时，$\gcd(F_m,F_n)=1$，所以$p_m\ne p_n$；否则同一素数将同时整除这两个互素数。由此构造出无穷多个不同素数。
\end{proof}
\paragraph{复习提示}不需要断言每个$F_n$都是素数，只需每个大于$1$的整数有素因子。

\subsection{习题14：莫比乌斯函数的约数和}\label{exercise:0:14}
\paragraph{题目}在正整数上定义$\mu(1)=1$；若$n$有素数平方因子，则$\mu(n)=0$；若$n=p_1\cdots p_r$是$r$个不同素数之积，则$\mu(n)=(-1)^r$。证明对$n>1$有
\[
\sum_{d\mid n}\mu(d)=0.
\]
\paragraph{考点与思路}只有无平方因子的约数有贡献；这些约数恰与不同素因子集合的子集一一对应。\textbf{原文依据：}本题的莫比乌斯函数定义，以及第零章第3节算术基本定理，书页11；原书本题未附证明，以下按素因数分解与定义补证。
\paragraph{所用结论}\toolref{n-prime-factor}、\toolref{n-binomial-count}。
\begin{proof}[详细证明]
由唯一分解，写$n=p_1^{\alpha_1}\cdots p_r^{\alpha_r}$，其中$p_i$两两不同，$\alpha_i\geq1$。由于$n>1$，$r\geq1$。

若$d\mid n$且$\mu(d)\ne0$，则$d$中每个素因子的指数只能是$0$或$1$，故存在唯一子集$J\subseteq\{1,\ldots,r\}$使
\[
d=\prod_{j\in J}p_j,\qquad \mu(d)=(-1)^{|J|}.
\]
空子集对应$d=1$，其贡献为$1$；其他含平方因子的约数贡献为$0$。按$|J|=k$分组，因此
\begin{align*}
\sum_{d\mid n}\mu(d)
&=\sum_{J\subseteq\{1,\ldots,r\}}(-1)^{|J|}\\
&=\sum_{k=0}^r\binom rk(-1)^k=(1-1)^r=0.
\end{align*}
\end{proof}
\paragraph{复习提示}条件$n>1$保证至少有一个素因子。当$n=1$时约数和为$\mu(1)=1$。

\subsection{习题23：解线性同余方程}\label{exercise:0:23}
\paragraph{题目}解$27x\equiv25\pmod{31}$。
\paragraph{考点与思路}先用欧几里得算法找逆元，再相乘。模$m$同余指$m$整除两数之差。\textbf{原文依据：}第零章第4节定理4，书页13；原证明先由互素性构造$au+nv=1$，将其乘以右端常数得到解，再用同余消去说明唯一性。以下保留此路线，具体以回代算出$u$。
\paragraph{所用结论}\toolref{n-gcd}、\toolref{n-bezout}、\toolref{n-congruence}。
\begin{proof}[详细解答]
由$31=27+4$、$27=6\cdot4+3$、$4=3+1$，回代得
\begin{align*}
1&=4-3=4-(27-6\cdot4)\\
 &=7\cdot4-27=7\cdot31-8\cdot27.
\end{align*}
所以$27$在模$31$下的逆元为$-8\equiv23$。原同余等价于
\[
x\equiv23\cdot25=575\equiv17\pmod{31}.
\]
检验：$27\cdot17-25=434=14\cdot31$。因$27$可逆，这一剩余类是唯一解。全部整数解为$x=17+31t$，$t\in\Z$。
\end{proof}
\paragraph{复习提示}同余方程的“唯一解”指模$31$唯一剩余类，整数解仍有无穷多个。

\subsection{习题24：解同余方程组}\label{exercise:0:24}
\paragraph{题目}解
\[
\begin{cases}
x\equiv2\pmod5,\\x\equiv1\pmod7,\\x\equiv5\pmod{11}.
\end{cases}
\]
\paragraph{考点与思路}沿用原文孙子定理证明中的构造：先求只在一个模数下同余于$1$、在其余模数下同余于$0$的整数，再按给定余数组合。\textbf{原文依据：}第零章第4节定理5（孙子定理），书页14；本题未附原书证明，以下按该定理的构造补写。
\paragraph{所用结论}\toolref{n-crt}、\toolref{n-bezout}、\toolref{n-congruence}。
\begin{proof}[详细解答]
模数$5,7,11$两两互素，令$M=5\cdot7\cdot11=385$，并令
\[
M_1=77,\qquad M_2=55,\qquad M_3=35.
\]
它们分别被另外两个模数整除。由于
\[
3\cdot77\equiv1\pmod5,\qquad
6\cdot55\equiv1\pmod7,\qquad
6\cdot35\equiv1\pmod{11},
\]
取$e_1=231$、$e_2=330$、$e_3=210$。则$e_i$在第$i$个模数下同余于$1$，在其余两个模数下同余于$0$。

如同原文定理5证明中以$x=b_1e_1+b_2e_2$构造两个条件的公共解，这里取三个对应项之和：
\begin{align*}
x_0&=2e_1+e_2+5e_3\\
&=2\cdot231+330+5\cdot210\\
&=1842=302+4\cdot385.
\end{align*}
对模$5$、模$7$、模$11$分别取余时，只有相应的那一项留下，因此$x_0$分别同余于$2,1,5$。故$302$也是一个解。

最后按原文证明的唯一性步骤：若$x,y$都是解，则$5,7,11$均整除$x-y$；它们两两互素，所以$385\mid x-y$。反过来，给任一解加$385k$不会改变三个余数。因此全部解为
\[
\boxed{x=302+385k,\qquad k\in\Z.}
\]
直接检验：$302-2=300$被$5$整除，$302-1=301$被$7$整除，$302-5=297$被$11$整除。
\end{proof}
\paragraph{复习提示}原文先处理两个模数，再说明可用归纳法推广；这里将同一个构造用于三个模数。给出一个解后，还须用互素性说明模$385$唯一。

\clearpage% !TEX root = main.tex
\section{第一章：代数结构}
本章选取作业中的19题。题目及证明思路依据现有笔记的习题文件和所引用的正文整理；下列细节是沿用这些思路的补充推导，不标作教材原证。各题“原文依据”指实际核对的现有笔记，未核对到教材证明的部分不声称已与教材原证核对。

\paragraph{本章常用记号与判别方法}
群的单位元记为$e$。子群$H\leq G$的左、右陪集分别为$aH$、$Ha$；左陪集个数记为指数$[G:H]$。若所有$a\in G$均满足$aH=Ha$，则称$H$为正规子群，记为$H\trianglelefteq G$。群中有消去律，但一般不能交换因子的次序。
非空子集$H$为子群的常用判别是：任意$x,y\in H$均有$xy^{-1}\in H$。环的双边理想$I$是加法子群，并且满足$RI\subseteq I$、$IR\subseteq I$。下文的“理想”均指双边理想。幂零元是某个正整数次幂为零的元素；零因子按笔记约定必须是非零元素。

\subsection{习题1：平方公式判别交换群}
\paragraph{题目}
若群$G$中任意$a,b$满足$(ab)^2=a^2b^2$，证明$G$交换。
\paragraph{考点与思路}
将幂展开，只消去两端的因子。也可证明该条件与交换性等价。
\textbf{原文依据：}现有笔记本题解答；群的消去律。
\paragraph{所用结论}
\toolref{g-laws}。
\begin{proof}[详细证明]
任取$a,b\in G$，题设给出$abab=aabb$。先在两端左乘$a^{-1}$，再右乘$b^{-1}$，得到
\begin{align*}
abab&=aabb,\\
bab&=abb,\\
ba&=ab.
\end{align*}
因为$a,b$任意，群中任意两元素交换，故$G$为交换群。

反过来，若$G$交换，则$abab=aabb=a^2b^2$，所以题设条件也成立。
\end{proof}
\paragraph{复习提示}
消去律仅允许消去位于同一侧的相同因子；不能从一开始就把中间的$ba$换成$ab$。

\subsection{习题2：每个元素的平方为单位元}
\paragraph{题目}
若群$G$中的每个元素$a$都满足$a^2=e$，证明$G$交换。
\paragraph{考点与思路}
将题设同时用于$a$、$b$和$ab$，再使用乘积取逆时的反序公式。
\textbf{原文依据：}现有笔记本题解答；乘积取逆公式。
\paragraph{所用结论}
\toolref{g-laws}。
\begin{proof}[详细证明]
由$a^2=e$知$a$本身就是$a$的逆元，所以每个$a\in G$都有$a^{-1}=a$。
任取$a,b\in G$，因为$ab\in G$，它也满足$(ab)^{-1}=ab$。另一方面，任何群中都有$(ab)^{-1}=b^{-1}a^{-1}$，于是
\[
ab=(ab)^{-1}=b^{-1}a^{-1}=ba.
\]
因此$G$交换。
\end{proof}
\paragraph{复习提示}
题设包含单位元，并不意味着每个元素的阶都恰为2；单位元的阶为1。结论也不要求$G$有限。

\subsection{习题4：可解方程的结合系统}
\paragraph{题目}
非空集合$G$上有一个结合的二元运算。若对任意$a,b\in G$，方程$ax=b$与$ya=b$都有解，证明$G$构成群。
\paragraph{考点与思路}
尚未证明是群时，不能使用逆元和消去律。先将某个元素的单侧单位元推广为全体元素的单侧单位元。
\textbf{原文依据：}现有笔记本题解答中先构造双边单位元、再构造双边逆元的步骤。
\paragraph{所用结论}
\toolref{g-group-construction}。
\begin{proof}[详细证明]
\textbf{第一步：构造双边单位元。}
因为$G$非空，可固定$a\in G$。由题设，存在$e,f\in G$使
\[
ae=a,\qquad fa=a.
\]
任取$x\in G$。方程$ya=x$有解，因此$x=ya$，结合律给出
\[
xe=(ya)e=y(ae)=ya=x.
\]
所以$e$是整个集合的右单位元。又由方程$az=x$有解，可写$x=az$，从而
\[
fx=f(az)=(fa)z=az=x.
\]
所以$f$是整个集合的左单位元。将两者用于$fe$，分别得到$fe=f$与$fe=e$，故$f=e$。因此$e$为双边单位元。

\textbf{第二步：构造双边逆元。}
任取$g\in G$，由方程可解性，存在$b,c\in G$使$gb=e$、$cg=e$。于是
\[
c=ce=c(gb)=(cg)b=eb=b.
\]
所以$gb=bg=e$，即$g$有双边逆元。

二元运算已经具有封闭性，结合律由题设保证；现在单位元和每个元素的逆元也已得到，故$G$是群。
\end{proof}
\paragraph{复习提示}
“有解”不必加强为“唯一解”。在构造出群以前，任何一步都不应预先使用消去律。

\subsection{习题6：任意偶数阶的非交换例子}
\paragraph{题目}
对整数$n>2$，构造一个$2n$阶的非交换群。
\paragraph{考点与思路}
沿用原解答，取正$n$边形的几何对称群，再按正文中相邻顶点像的计数和旋转、反射关系证明其阶数与非交换性。
\textbf{原文依据：}现有笔记本题解答；正文“图形的对称群与二面体群”中对称群命题、二面体群定义及阶数与非交换性的命题。
\paragraph{所用结论}
\toolref{g-dihedral}。
\begin{proof}[详细证明]
取平面上的正$n$边形，记$D_n$为保持此正多边形整体不变的所有平面等距变换的集合。恒等变换属于$D_n$；两个对称变换的复合仍保持图形；逆变换也保持图形。映射复合具有结合律，所以$D_n$在复合下成群。

\textbf{先证明$|D_n|=2n$。}任意对称变换置换顶点，因而保持顶点的重心，即正多边形的中心$O$。固定两个相邻顶点$P,Q$，由于$n>2$，$O,P,Q$不共线。平面等距变换由三个不共线点的像唯一决定，因此一个对称变换由$P,Q$的像唯一决定。
$P$的像有$n$种选择；$Q$的像必须与$P$的像相邻，而每个顶点恰有两个相邻顶点，所以$Q$的像至多有两种选择。故$|D_n|\leq2n$。

令$T$为绕$O$旋转$2\pi/n$，$S$为关于一条固定对称轴的反射。变换
\[
T^k,\quad ST^k\qquad(0\leq k<n)
\]
都保持正多边形。$n$个旋转$T^k$的旋转角在模$2\pi$意义下不同，故彼此不同；左复合$S$可逆，故$n$个$ST^k$也彼此不同。旋转保持平面定向，反射与旋转的复合反转平面定向，所以两类不相交。已经得到$2n$个不同对称变换，结合前面的上界，得$|D_n|=2n$。

\textbf{再证明它不交换。}反射反转旋转方向，所以$STS=T^{-1}$，且$S^2=e$。若$ST=TS$，则右乘$S$得$STS=T$，与前式比较可得$T=T^{-1}$，即$T^2=e$。
但是$T$的阶恰为$n$：旋转$k$次恢复原位置当且仅当$2k\pi/n$为$2\pi$的整数倍，最小正整数$k$为$n$。所以$T^2=e$将给出$n\mid2$，与$n>2$矛盾。故$ST\ne TS$，$D_n$为所需的非交换群。
\end{proof}
\paragraph{复习提示}
本笔记约定$|D_n|=2n$。条件$n>2$不可略去，2阶和4阶群均交换。

\subsection{习题7：反复共轭}
\paragraph{题目}
若群$G$中的$a,b$满足$a^{-1}ba=b^r$，其中$r$为正整数，证明对任意正整数$i$都有$a^{-i}ba^i=b^{r^i}$。
\paragraph{考点与思路}
先验证共轭保持正整数次幂，再对共轭次数$i$归纳。
\textbf{原文依据：}现有笔记本题解答中的共轭幂公式与数学归纳法。
\paragraph{所用结论}
\toolref{g-laws}、\toolref{n-induction}。
\begin{proof}[详细证明]
对任意正整数$m$，将相邻的$aa^{-1}$消去，得
\[
(a^{-1}ba)^m=a^{-1}b^ma.
\]
因此$a^{-1}b^ma=(b^r)^m=b^{rm}$。

当$i=1$时，结论就是题设。假设某个$i\geq1$已有$a^{-i}ba^i=b^{r^i}$，则
\begin{align*}
a^{-(i+1)}ba^{i+1}
&=a^{-1}(a^{-i}ba^i)a\\
&=a^{-1}b^{r^i}a\\
&=b^{r\cdot r^i}=b^{r^{i+1}}.
\end{align*}
归纳法给出所有正整数$i$的结论。
\end{proof}
\paragraph{复习提示}
反复共轭使指数每次乘以$r$，最终指数为$r^i$，不是$ri$。

\subsection{习题8：取逆映射的同态性}
\paragraph{题目}
证明：群$G$交换，当且仅当映射$\iota:G\to G$、$\iota(a)=a^{-1}$是群同态。取逆映射本身为双射，所以本题中的同态条件也等价于它是同构。
\paragraph{考点与思路}
取逆映射总为双射；真正需要判别的是它是否保乘法。
\textbf{原文依据：}现有笔记本题解答中先验证取逆为双射、再比较保乘法与反序公式的步骤。
\paragraph{所用结论}
\toolref{g-laws}、\toolref{g-first-iso}。
\begin{proof}[详细证明]
因为$\iota(\iota(a))=(a^{-1})^{-1}=a$，所以$\iota$的逆映射就是自身，因而总是双射。

\textbf{先由同态推出交换性。}若$\iota$为同态，则保乘法给出
\[
a^{-1}b^{-1}=\iota(a)\iota(b)=\iota(ab)=b^{-1}a^{-1}.
\]
两边取逆得到$ba=ab$。因$a,b$任意，$G$交换。

\textbf{再由交换性推出同构。}若$G$交换，则对任意$a,b\in G$，
\[
\iota(ab)=(ab)^{-1}=b^{-1}a^{-1}=a^{-1}b^{-1}=\iota(a)\iota(b).
\]
所以$\iota$为同态，再结合双射性，知它是同构。
\end{proof}
\paragraph{复习提示}
任何群都满足$(ab)^{-1}=b^{-1}a^{-1}$；保乘法所要求的却是$(ab)^{-1}=a^{-1}b^{-1}$。

\subsection{习题9：由子集定义等价关系}
\paragraph{题目}
设$S\subseteq G$非空，在$G$上定义$a\sim b\iff ab^{-1}\in S$。证明该关系为等价关系，当且仅当$S$为子群。
\paragraph{考点与思路}
等价关系要求自反、对称、传递。反向将关系与单位元联系起来，导出子群判别条件。
\textbf{原文依据：}现有笔记本题解答；子群判别法与右陪集。
\paragraph{所用结论}
\toolref{g-subgroup}、\toolref{g-cosets}。
\begin{proof}[详细证明]
\textbf{若$S$为子群。}任取$a\in G$，因$aa^{-1}=e\in S$，有$a\sim a$，满足自反性。
若$a\sim b$，则$ab^{-1}\in S$，子群对逆封闭给出$ba^{-1}=(ab^{-1})^{-1}\in S$，所以$b\sim a$。
若$a\sim b$且$b\sim c$，则$ab^{-1},bc^{-1}\in S$，因而
\[
ac^{-1}=(ab^{-1})(bc^{-1})\in S,
\]
即$a\sim c$，满足传递性。因此该关系是等价关系。

\textbf{若该关系为等价关系。}自反性给出$e=aa^{-1}\in S$。任取$u,v\in S$，有$u\sim e$、$v\sim e$。由对称性得$e\sim v$，再由传递性得$u\sim v$，故$uv^{-1}\in S$。

说明此条件为何足够：取$u=e$可知$v^{-1}\in S$；再将已经得到的$v^{-1}$代入上述差式条件，得到$u(v^{-1})^{-1}=uv\in S$。所以$S$含单位元并对逆和乘积封闭，确为子群。

此外，$a\sim b$等价于$a=sb$对某个$s\in S$成立，因此$b$所在的等价类为右陪集$Sb$。
\end{proof}
\paragraph{复习提示}
$ab^{-1}$对应右陪集$Sb$；不要与$a^{-1}b$所对应的左陪集混淆。

\subsection{习题10：整数加法群与其倍数子群}
\paragraph{题目}
设$n$为正整数，证明加法群$n\Z$与$\Z$同构。
\paragraph{考点与思路}
同构必须写出映射，并分别验证保运算、单射和满射。
\textbf{原文依据：}现有笔记本题解答中的映射$k\mapsto nk$。
\paragraph{所用结论}
\toolref{g-first-iso}、\toolref{r-quotients}（其中$\Z$为整环，非零整数可消去）。
\begin{proof}[详细证明]
定义$\varphi:\Z\to n\Z$为$\varphi(k)=nk$。对任意$k,l\in\Z$，
\[
\varphi(k+l)=n(k+l)=nk+nl=\varphi(k)+\varphi(l),
\]
所以它保加法。若$\varphi(k)=\varphi(l)$，则$n(k-l)=0$。因为$n\ne0$，整数乘法无零因子，故$k=l$，所以它为单射。
任取$n\Z$的元素，按定义可写成$nk$，它就是$\varphi(k)$，所以映射满射。因此$\Z\cong n\Z$。
\end{proof}
\paragraph{复习提示}
这里讨论加法群。若$n>1$，上述映射不保通常的乘法，例如$\varphi(1\cdot1)=n$，而$\varphi(1)\varphi(1)=n^2$。若允许$n=0$，同构结论不成立。

\subsection{习题12：指数2的子群正规}
\paragraph{题目}
若$G$为$2n$阶群，其中有一个$n$阶子群$H$，证明$H$正规。原题的条件给出$[G:H]=2$；以下按现有笔记的左右陪集方法证明这一更一般的指数2结论。
\paragraph{考点与思路}
指数2意味着除$H$以外，只有一个左陪集和一个右陪集，它们都必须是$H$的补集。
\textbf{原文依据：}现有笔记本题解答中按$a\in H$与$a\notin H$分类、比较左右陪集的证明。
\paragraph{所用结论}
\toolref{g-cosets}、\toolref{g-normal}、\toolref{g-lagrange}。
\begin{proof}[详细证明]
任取$a\in G$。若$a\in H$，子群性质给出$aH=H=Ha$。

若$a\notin H$，则$aH\ne H$，并且不同左陪集不交。因为恰有两个左陪集，必有$aH=G\setminus H$。
取逆映射将左陪集$xH$一一对应到右陪集$Hx^{-1}$，所以右陪集也恰有两个。同理$Ha=G\setminus H$。于是这一情形下也有$aH=Ha$。

对所有$a\in G$均已证明$aH=Ha$，故$H$正规。有限群的特殊情形由$[G:H]=|G|/|H|=2$立即得到。
\end{proof}
\paragraph{复习提示}
一般结论不要求群有限；只有用阶数比计算指数时才使用有限性。

\subsection{习题18：只有有限多个子群的群必有限}
\paragraph{题目}
证明：如果群$G$只有有限多个子群，那么$G$必为有限群。
\paragraph{考点与思路}
先排除无限阶元素，再将全群表示为有限多个有限循环子群的并。
\textbf{原文依据：}现有笔记本题解答中先排除无限阶元素、再取有限多个循环子群之并的证明。
\paragraph{所用结论}
\toolref{g-cyclic}。
\begin{proof}[详细证明]
\textbf{第一步：每个元素都为有限阶。}
反设$a\in G$为无限阶，则子群$\langle a^m\rangle$（$m=1,2,\ldots$）两两不同。
事实上，若$\langle a^m\rangle=\langle a^n\rangle$，则存在整数$u,v$使$a^m=a^{nu}$、$a^n=a^{mv}$。无限阶意味着相等幂的指数相等，故$m=nu$、$n=mv$。由于$m,n$为正整数，$u,v$均为正整数；于是$uv=1$，得$u=v=1$、$m=n$。
这产生无限多个子群，与假设矛盾。

\textbf{第二步：全群有限。}
所有循环子群均为$G$的子群，所以它们只有有限多个，记为$C_1,\ldots,C_t$。第一步说明每个$C_j$都有限。
任意$g\in G$属于其生成的循环子群$\langle g\rangle$，因此
\[
G=\bigcup_{j=1}^{t}C_j,\qquad |G|\leq\sum_{j=1}^{t}|C_j|<\infty.
\]
所以$G$有限。
\end{proof}
\paragraph{复习提示}
仅证明“每个元素有限阶”还不足以证明群有限；必须利用循环子群的个数也有限。

\subsection{习题21：有限子群乘积的计数}
\paragraph{题目}
设$H,K$为群$G$的有限子群，证明
\[
|HK|=\frac{|H||K|}{|H\cap K|}.
\]
此处不要求集合$HK$是子群。
\paragraph{考点与思路}
映射$(h,k)\mapsto hk$一般不单射；准确计算每个像的原像个数。
\textbf{原文依据：}现有笔记本题解答中的满射$H\times K\to HK$及每个像的原像计数。
\paragraph{所用结论}
\toolref{g-laws}、\toolref{g-fiber-count}。
\begin{proof}[详细证明]
定义满射$\pi:H\times K\to HK$，$\pi(h,k)=hk$。固定$g=hk\in HK$的一种表示。

若$h'k'=hk$，两边左乘$h^{-1}$、右乘$k'^{-1}$，得到
\[
t:=h^{-1}h'=kk'^{-1}\in H\cap K.
\]
由此$h'=ht$、$k'=t^{-1}k$。反之，任取$t\in H\cap K$，都有$ht\in H$、$t^{-1}k\in K$，且
\[
(ht)(t^{-1}k)=hk=g.
\]
不同$t$给出不同的有序对，因为$ht_1=ht_2$可消去$h$而得到$t_1=t_2$。因此$g$的全部原像与$H\cap K$一一对应，恰有$|H\cap K|$个。

对有限集合$H\times K$按原像分组计数，得
\[
|H||K|=|HK|\,|H\cap K|,
\]
除以非零正整数$|H\cap K|$即得所求。
\end{proof}
\paragraph{复习提示}
需要扣除的是交集造成的重复表示；即使$HK$不是子群，集合计数公式仍成立。

\subsection{习题22：正规子群的乘积}
\paragraph{题目}
设$M,N\trianglelefteq G$，证明$MN=NM$且$MN\trianglelefteq G$。若$M\cap N=\{e\}$，证明$MN/N\cong M$。
\paragraph{考点与思路}
先证明集合乘积是正规子群，再构造商群的满同态并求核。
\textbf{原文依据：}现有笔记本题解答中先证乘积相等与正规、再构造$m\mapsto mN$的步骤；同态基本定理。
\paragraph{所用结论}
\toolref{g-subgroup}、\toolref{g-normal}、\toolref{g-first-iso}。
\begin{proof}[详细证明]
\textbf{第一步：证明$MN=NM$。}
任取$m\in M,n\in N$，因$M$正规，$n^{-1}mn\in M$，于是$mn=n(n^{-1}mn)\in NM$，故$MN\subseteq NM$。
同样，由$M$正规，$nmn^{-1}\in M$，所以$nm=(nmn^{-1})n\in MN$，得到反向包含。

\textbf{第二步：证明$MN$为正规子群。}
$e\in MN$。集合乘积满足结合律，利用$NM=MN$，得
\[
(MN)(MN)=M(NM)N=M(MN)N=MN,
\qquad (MN)^{-1}=NM=MN.
\]
所以$MN$对乘法和逆封闭，为子群。再任取$g\in G$，由两个子群的正规性，
\[
g(MN)g^{-1}=(gMg^{-1})(gNg^{-1})=MN.
\]
因此$MN\trianglelefteq G$。

\textbf{第三步：构造同构。}
因为$N$在$G$中正规，它在子群$MN$中也正规，可以构成商群$MN/N$。定义
\[
\varphi:M\longrightarrow MN/N,\qquad \varphi(m)=mN.
\]
商群运算给出$\varphi(m_1m_2)=m_1m_2N=(m_1N)(m_2N)$，所以$\varphi$是同态。
任意陪集$mnN$等于$mN$，故$\varphi$满射。其核为
\[
\ker\varphi=\{m\in M:mN=N\}=M\cap N.
\]
因此同态基本定理给出$M/(M\cap N)\cong MN/N$。当$M\cap N=\{e\}$时，$\varphi$本身单射，故$M\cong MN/N$。
\end{proof}
\paragraph{复习提示}
证明商群同构时，应逐项交代映射、保运算、满射和核；交平凡正是单射的依据。

\subsection{习题29：改变坐标得到同构环}
\paragraph{题目}
在交换幺环$R$的底集上定义
\[
a\oplus b=a+b-1,\qquad a\odot b=a+b-ab.
\]
证明新运算使它成为与原环同构的交换幺环，并求新的零元、单位元及加法逆元。
\paragraph{考点与思路}
寻找坐标变换$f(a)=1-a$，把新运算转回熟悉的运算。
\textbf{原文依据：}现有笔记本题解答中的坐标变换$f(a)=1-a$。
\paragraph{所用结论}
\toolref{g-transport}。
\begin{proof}[详细证明]
记具有新运算的集合为$R_{\mathrm{new}}$，定义$f:R_{\mathrm{new}}\to R$为$f(a)=1-a$。因为$f(f(a))=1-(1-a)=a$，它是双射。
对任意$a,b\in R$，
\begin{align*}
f(a\oplus b)&=1-(a+b-1)=2-a-b=f(a)+f(b),\\
f(a\odot b)&=1-(a+b-ab)=1-a-b+ab\\
&=(1-a)(1-b)=f(a)f(b).
\end{align*}
所以$f$将两种新运算分别对应到原环的加法、乘法。由于$f$单射，原环的每一条运算恒等式均可通过$f$拉回；例如
\begin{align*}
f((a\odot b)\odot c)&=(f(a)f(b))f(c)\\
&=f(a)(f(b)f(c))=f(a\odot(b\odot c)),
\end{align*}
从而新乘法满足结合律。加法的结合、交换以及分配律同理由原环的相应公理得到。因此新结构为环，且$f$为环同构。

新零元是$f^{-1}(0)=1$，新乘法单位元是$f^{-1}(1)=0$。也可直接核对$a\oplus1=a$、$a\odot0=a$。元素$a$的新加法逆元为$2-a$，因为
\[
a\oplus(2-a)=a+(2-a)-1=1,
\]
右端正是新零元。新乘法交换性由$a+b-ab=b+a-ba$得到。
\end{proof}
\paragraph{复习提示}
式中$0,1,2$表示原环元素，$2=1+1$；新结构的零元与单位元分别是原来的$1$与$0$。

\subsection{习题30：子环与母环的单位元和交换性}
\paragraph{题目}
分别举例说明：(i)母环有单位元而子环无单位元；(ii)母环无单位元而子环有单位元；(iii)母环和子环都有单位元，但二者不同；(iv)母环不交换而子环交换。
\paragraph{考点与思路}
这里子环只要求对加法、取负和乘法封闭，并不要求与母环共用单位元。应检查实例确实满足对应性质。
\textbf{原文依据：}现有笔记本题解答的四个反例构造及原有顺序。
\paragraph{所用结论}
\toolref{g-ring-criteria}。
\begin{proof}[详细证明]
\textbf{(i)}取$R=\Z$、$S=2\Z$。偶数对差和积封闭，所以$S$是子环；$R$的单位元为1。若$u\in2\Z$是$S$的单位元，则$u\cdot2=2$，在整数中得$u=1$，与$u$为偶数矛盾。因此$S$没有单位元。

\textbf{(ii)}取逐分量运算的环$R=2\Z\times\Z$，以及子环$S=\{0\}\times\Z$。若$R$有单位元$(u,v)$，将它乘以$(2,0)$得到$(2u,0)=(2,0)$，所以$u=1\notin2\Z$，矛盾。而对任意$(0,k)\in S$，
\[
(0,1)(0,k)=(0,k)=(0,k)(0,1),
\]
所以$S$有单位元$(0,1)$。

\textbf{(iii)}取$R=M_2(\R)$，令
\[
S=\left\{\begin{pmatrix}a&0\\0&0\end{pmatrix}:a\in\R\right\}.
\]
两个这样的矩阵之差和积仍有相同形状，所以$S$是子环。母环的单位元为$E=\operatorname{diag}(1,1)$，子环的单位元为$\operatorname{diag}(1,0)$，两者不同。

\textbf{(iv)}仍取$R=M_2(\R)$，令$S=\{aE:a\in\R\}$为数量矩阵构成的子环。因$(aE)(bE)=abE=baE$，$S$交换。记$E_{ij}$为第$(i,j)$项为1、其余项为0的矩阵，则
\[
E_{12}E_{21}=E_{11},\qquad E_{21}E_{12}=E_{22},
\]
二者不同，故母环不交换。
\end{proof}
\paragraph{复习提示}
单位元是相对于环内全部元素定义的；子环中的单位元不必对母环其他元素起单位作用。

\subsection{习题32：域没有非平凡双边理想}
\paragraph{题目}
证明域$F$的双边理想只有$\{0\}$和$F$。同一结论也适用于体。
\paragraph{考点与思路}
非零理想中取非零元素，用其逆元推出单位元属于该理想。
\textbf{原文依据：}现有笔记本题解答中由非零理想元素推出$1\in I$的证明。
\paragraph{所用结论}
\toolref{g-ring-criteria}。
\begin{proof}[详细证明]
设$I$为$F$的理想。若$I=\{0\}$，已经属于所列情形。若$I\ne\{0\}$，可取$a\in I$、$a\ne0$。域中$a$有逆元$a^{-1}\in F$，理想的吸收性给出$1=a^{-1}a\in I$。
任取$r\in F$，再次由吸收性，$r=r\cdot1\in I$，所以$F\subseteq I$。结合$I\subseteq F$，得到$I=F$。

证明只使用非零元素有双边逆元，没有使用乘法交换性，所以对不一定交换的体也成立。
\end{proof}
\paragraph{复习提示}
一般环中非零元素未必可逆，不能对任意环照搬这一步。

\subsection{习题34：理想的根仍为理想}
\paragraph{题目}
设$R$为交换环、$I$为其理想，定义
\[
\sqrt I=\{r\in R:\text{存在正整数 }n\text{ 使 }r^n\in I\}.
\]
证明$\sqrt I$为$R$的理想。
\paragraph{考点与思路}
验证差封闭和吸收性。两个元素对应的幂次数可能不同，统一选取$n+m-1$次幂。
\textbf{原文依据：}现有笔记本题解答中的$(r-s)^{n+m-1}$二项式展开和吸收性验证。
\paragraph{所用结论}
\toolref{g-ring-criteria}、\toolref{g-binomial}。
\begin{proof}[详细证明]
首先$0^1=0\in I$，故$0\in\sqrt I$，集合非空。

任取$r,s\in\sqrt I$，由定义可选正整数$n,m$使$r^n\in I$、$s^m\in I$。令$N=n+m-1$。因为环交换，可以使用二项式展开：
\[
(r-s)^N=\sum_{k=0}^{N}(-1)^{N-k}\binom Nk r^k s^{N-k}.
\]
对任一项，若$k\geq n$，它含有因子$r^n$，故属于$I$；若$k<n$，则$k\leq n-1$，从而$N-k\geq m$，这一项含有因子$s^m$，也属于$I$。整数倍与有限和仍在加法子群$I$中，故$(r-s)^N\in I$，即$r-s\in\sqrt I$。因此$\sqrt I$为加法子群。

任取$t\in R$、$r\in\sqrt I$，选$n\geq1$使$r^n\in I$。由交换性和理想吸收性，
\[
(tr)^n=t^nr^n\in I.
\]
所以$tr\in\sqrt I$；交换性又给出$rt=tr\in\sqrt I$。因此$\sqrt I$满足双边吸收性，确为理想。

若$R$没有单位元，上述展开仍然成立：在$k=0$或$k=N$的端项中，直接写作$(-s)^N$或$r^N$，不引入$R$内不存在的零次幂。判断因子属于$I$时，若剩余因子个数为零，就直接使用$r^n\in I$或$s^m\in I$。
\end{proof}
\paragraph{复习提示}
不能仅由$r^n,s^m\in I$断言$(r-s)^{\max(n,m)}\in I$；$n+m-1$保证展开中每一项都含足够多的某一种因子。

\subsection{习题39：交换环的幂零元构成理想}
\paragraph{题目}
证明交换环中全体幂零元构成理想，并说明非交换环中该结论不成立。
\paragraph{考点与思路}
它是习题34中$I=\{0\}$的特例；复习时也应能独立写出幂次数的论证。
\textbf{原文依据：}现有笔记本题解答；习题34取$I=\{0\}$的情形，以及矩阵单位$E_{12},E_{21}$的非交换反例。
\paragraph{所用结论}
\toolref{g-ring-criteria}、\toolref{g-binomial}。
\begin{proof}[详细证明]
记$\mathcal N=\{a\in R:\text{存在正整数 }n\text{ 使 }a^n=0\}$。因为$0\in\mathcal N$，它非空。
若$a,b\in\mathcal N$，选$n,m\geq1$使$a^n=b^m=0$。展开$(a-b)^{n+m-1}$，其中任意一项$a^k b^{n+m-1-k}$都满足$k\geq n$或$n+m-1-k\geq m$，因而为零。所以
\[
(a-b)^{n+m-1}=0,
\]
即$a-b\in\mathcal N$。若$r\in R$，交换性给出$(ra)^n=r^na^n=0$，所以$ra=ar\in\mathcal N$。因此$\mathcal N$为理想。

非交换环的反例取$M_2(\Q)$中的
\[
A=\begin{pmatrix}0&1\\0&0\end{pmatrix},\qquad
B=\begin{pmatrix}0&0\\1&0\end{pmatrix}.
\]
有$A^2=B^2=0$，但$(A+B)^2=E$。若$A+B$幂零，某个正整数$k$应满足$(A+B)^k=0$，于是$(A+B)^{2k}=0$；另一方面$(A+B)^{2k}=E$，矛盾。因此幂零元之和未必幂零，幂零元集甚至不一定对加法封闭。
\end{proof}
\paragraph{复习提示}
证明中交换性的用途是二项式展开以及$(ra)^n=r^na^n$，不能把这两步用于任意非交换环。

\subsection{习题43：连续函数环中的理想与零因子}
\paragraph{题目}
记$C[0,1]$为$[0,1]$上全部实值连续函数构成的环，加法和乘法逐点定义。
\begin{enumerate}
\item 每个真理想$I\ne C[0,1]$都存在公共零点，即存在$t_0\in[0,1]$使所有$f\in I$均满足$f(t_0)=0$。
\item 对非零函数$f\in C[0,1]$，证明$f$为零因子，当且仅当其零点集包含一个非空的相对开区间。
\end{enumerate}
相对开区间指$(\alpha,\beta)\cap[0,1]$形式的非空集合，其中$\alpha<\beta$为实数。这里“非零函数”指不恒为零的函数。
\paragraph{考点与思路}
第一问用紧致性选有限子覆盖，再以平方和构造处处非零的理想元素；第二问使用连续函数在非零点附近仍非零的性质。
\textbf{原文依据：}现有笔记本题解答中紧致性、有限平方和的构造，以及非零邻域和三角形函数的证明。
\paragraph{所用结论}
\toolref{g-ring-criteria}、\toolref{g-analysis}。
\begin{proof}[详细证明]
\textbf{(i) 真理想具有公共零点。}
反设$I$没有公共零点，则对每个$t\in[0,1]$，都可选$f_t\in I$使$f_t(t)\ne0$。
由连续性存在$t$的相对开邻域$U_t$，使$f_t$在$U_t$上处处非零。这些邻域覆盖$[0,1]$。

闭区间$[0,1]$紧致，故可选有限子覆盖$U_{t_1},\ldots,U_{t_m}$。记$f_j=f_{t_j}$，并令
\[
g=f_1^2+\cdots+f_m^2.
\]
因为$f_j\in I$，理想吸收性给出$f_j^2\in I$，加法封闭性给出$g\in I$。
任意$x\in[0,1]$至少属于某个$U_{t_j}$，因此对应$f_j(x)\ne0$，从而$g(x)=\sum_{j=1}^m f_j(x)^2>0$。
所以函数$h(x)=1/g(x)$在整个区间连续，属于$C[0,1]$。

由理想吸收性，常值函数$1=gh\in I$。于是任取$u\in C[0,1]$都有$u=u\cdot1\in I$，这说明$I=C[0,1]$，与$I$为真理想矛盾。故必有公共零点。

\textbf{(ii) 零因子的必要性。}
设$f\ne0$为零因子，则存在不恒为零的$g\in C[0,1]$使$fg=0$。选$t\in[0,1]$满足$g(t)\ne0$。
连续性给出某个$\delta>0$，使得当$x\in[0,1]$且$|x-t|<\delta$时，
\[
|g(x)-g(t)|<\frac{|g(t)|}{2},\qquad
|g(x)|>\frac{|g(t)|}{2}>0.
\]
在非空相对开区间$(t-\delta,t+\delta)\cap[0,1]$上，$f(x)g(x)=0$且$g(x)\ne0$，所以$f(x)=0$。因此$f$的零点集包含这样的区间。

\textbf{(ii) 零因子的充分性。}
反过来，设$f$在非空相对开区间$J$上恒为零。这样的$J$含有一个位于$(0,1)$内的小闭区间$[c-\varepsilon,c+\varepsilon]$，其中$\varepsilon>0$。这是因为非空相对开集即使包含端点，也包含相邻的内部点，故可在其内部选$c$及足够小的$\varepsilon$。
定义连续的三角形函数
\[
g(x)=\max\left\{0,1-\frac{|x-c|}{\varepsilon}\right\},\qquad x\in[0,1].
\]
它满足$g(c)=1$，所以不恒为零；在$[c-\varepsilon,c+\varepsilon]$外为零，而在该闭区间内$f$为零。因此对所有$x\in[0,1]$均有$f(x)g(x)=0$，即$fg=0$。由于$f\ne0$，$f$为零因子。
\end{proof}
\paragraph{复习提示}
“每个$f\in I$各有零点”与“所有$f\in I$有同一个公共零点”的量词次序不同。第二问不能只要求$f$有一个零点；必须有一段区间。零函数满足零点集条件，但按约定不称零因子，所以题目须限定$f\ne0$。

\subsection{习题44：有限域上的矩阵计数}
\paragraph{题目}
设$p$为素数，$F=\F_p$为$p$元域，$n\geq1$。证明
\[
|M_n(F)|=p^{n^2},\qquad
|\GL_n(F)|=\prod_{k=0}^{n-1}(p^n-p^k).
\]
这里$M_n(F)$是全体$n$阶矩阵的集合，$\GL_n(F)$是其中的可逆矩阵乘法群。
\paragraph{考点与思路}
可逆矩阵的列向量恰为$F^n$的一组有序基。逐列选取时，排除前面列向量张成的子空间。
\textbf{原文依据：}现有笔记本题解答中按矩阵位置计数、再按线性无关列向量逐列计数的证明。
\paragraph{所用结论}
\toolref{g-basis-count}。
\begin{proof}[详细证明]
一个$n$阶矩阵共有$n^2$个位置，每一位置均可独立选取$F$的$p$个元素之一。因此乘法计数原理给出$|M_n(F)|=p^{n^2}$。

现在计数可逆矩阵。线性代数的可逆性判别说明，矩阵可逆当且仅当其$n$个列向量线性无关。第一列不能是零向量，所以有$p^n-1$种选择。

假设已选出线性无关的前$k$列$v_1,\ldots,v_k$，其中$0\leq k<n$。其张成空间中的向量均可写成
\[
\lambda_1v_1+\cdots+\lambda_kv_k,\qquad \lambda_i\in F.
\]
因为这$k$列线性无关，不同系数组给出不同向量：若两组给出相同向量，相减后得到线性关系，所有系数差必须为零。因此该张成空间恰有$p^k$个向量。
第$k+1$列必须位于该空间之外，才能保持列向量线性无关，所以有$p^n-p^k$种选择。

依次选满$n$列，每个有序选择恰对应一个矩阵，没有额外的排列除数。故
\[
|\GL_n(F)|=(p^n-1)(p^n-p)\cdots(p^n-p^{n-1}).
\]
这正是所求乘积。
\end{proof}
\paragraph{复习提示}
当$n=1$时结果为$p-1$，对应域中非零元素的个数。关键区别是：每一列都非零不足以保证可逆，还必须排除前面列向量的整个张成空间。

\clearpage% 根据现有 chapter2-exercises.tex 展开；本文件不声称与未复核的原书证明相同。
\section{第二章：群论}

\subsection*{记号与必备工具}
本章置换的乘积按从右到左的顺序作用。$S_n$为$n$个点的对称群，$A_n$为其中偶置换组成的交错群。对换是交换两个点的置换；每个置换可分解为对换的乘积。每个偶置换可分解为3-轮换的乘积：将对换分解两两配对，相同对换相消，共用一个点的两个对换之积为3-轮换，不相交时可用
\[
(ab)(cd)=(acb)(acd)
\]
分解。故证明某子群等于$A_n$，可设法生成全部3-轮换。

群的中心、子群的正规化子分别为
\[
Z(G)=\{z\in G:zg=gz\ (g\in G)\},\qquad
N_G(H)=\{g\in G:gHg^{-1}=H\}.
\]
记$[x,y]=x^{-1}y^{-1}xy$为换位子，$G^{(1)}=\langle[x,y]:x,y\in G\rangle$为换位子群（导群），递归定义$G^{(0)}=G$、$G^{(r+1)}=(G^{(r)})^{(1)}$。存在$r$使$G^{(r)}=\{e\}$时，称$G$可解。群交换当且仅当它的导群平凡。

若有限群阶为$p^am$，$p\nmid m$，则$p^a$阶子群称西罗$p$子群。西罗定理说明这样的子群存在，全部互相共轭，其个数$n_p$满足
\[
n_p\mid m,\qquad n_p\equiv1\pmod p,
\]
且$n_p=1$当且仅当该西罗子群正规。
共轭作用给出子群$H$的共轭子群数为$[G:N_G(H)]$。

自同构是群到自身的同构，全部自同构组成$\Aut(G)$。映射$c_g:x\mapsto gxg^{-1}$称内自同构，全部内自同构组成$\operatorname{Inn}(G)$。同态$g\mapsto c_g$的核为$Z(G)$，所以
\[
\operatorname{Inn}(G)\cong G/Z(G).
\]
中心平凡且每个自同构都是内自同构的群称为\textbf{完全群}。只有平凡正规子群和自身、且本身非平凡的群称为\textbf{单群}。

以下各题保留现有笔记的题号；证明为在现有解答思路上展开的补充证明。

\subsection{习题7：有理数加法群的有限生成子群}
\paragraph{题目}证明$(\Q,+)$的任一有限生成子群都是循环群。
\paragraph{考点与思路}把有限多个有理数通分，将生成子群转化为整数的线性组合，再用最大公因数和贝祖等式确定生成元。
\textbf{原文依据：}现有笔记本题解答（通分与贝祖等式）。
\paragraph{所用结论}\toolref{n-bezout}（整数线性组合）与\toolref{g-cyclic}（循环群与整数子群）。
\begin{proof}[详细证明]
设$H=\langle q_1,\ldots,q_t\rangle$，这里加法群中的生成意味着
\[
H=\left\{\sum_{i=1}^t k_iq_i:k_i\in\Z\right\}.
\]
若生成集为空或全部$q_i=0$，则$H=\{0\}=\langle0\rangle$，已经循环。以下设至少一个$q_i\ne0$。

取所有$q_i$的正公分母$D$，写$q_i=m_i/D$，其中$m_i\in\Z$。令$d=\gcd(m_1,\ldots,m_t)>0$。因$d$整除每个$m_i$，对任意$k_i\in\Z$都有
\[
\sum_i k_iq_i=\frac{\sum_i k_im_i}{D}\in\frac dD\Z.
\]
因此$H\subseteq(d/D)\Z$。

反过来，由贝祖等式，存在整数$u_1,\ldots,u_t$使$\sum_i u_im_i=d$。除以$D$得
\[
\frac dD=\sum_i u_iq_i\in H.
\]
由于$H$是加法子群，它包含$d/D$的全部整数倍，故$(d/D)\Z\subseteq H$。综上，
\[
H=\frac dD\Z=\left\langle\frac dD\right\rangle.
\]
\end{proof}
\paragraph{复习提示}“每个有限生成子群循环”不表示$\Q$本身循环。若$\Q=\langle q\rangle$且$q\ne0$，则$q/2$不是$q$的整数倍，矛盾。旧记号$\Q^+$在此指有理数加法群。

\subsection{习题8：有限阶生成元与交换性}
\paragraph{题目}交换群$G$由有限多个有限阶元素生成，证明$G$有限。
\paragraph{考点与思路}交换性允许集中同一生成元，有限阶允许把指数限制在有限范围内。
\textbf{原文依据：}现有笔记本题解答（交换重排与指数取余）。
\paragraph{所用结论}\toolref{g-laws}（交换重排与幂运算）及\toolref{g-cyclic}（有限阶的指数取余）。
\begin{proof}[详细证明]
设$G=\langle a_1,\ldots,a_t\rangle$，记$n_i=\operatorname{ord}(a_i)<\infty$。由生成的定义，任意$g\in G$是有限多个$a_i$和$a_i^{-1}$的乘积。由于$G$交换，可重排这些因子，将同类因子合并，得到
\[
g=a_1^{k_1}\cdots a_t^{k_t},\qquad k_i\in\Z.
\]
对每个指数作带余除法$k_i=q_in_i+r_i$，其中$0\le r_i<n_i$。由$a_i^{n_i}=e$，有$a_i^{k_i}=a_i^{r_i}$。故映射
\[
\prod_{i=1}^t\{0,1,\ldots,n_i-1\}\longrightarrow G,
\qquad(r_1,\ldots,r_t)\longmapsto a_1^{r_1}\cdots a_t^{r_t}
\]
满射。定义域只有$\prod_i n_i$个元素，因而
\[
|G|\le\prod_{i=1}^t n_i<\infty.
\]
空生成集对应平凡群，也满足结论。
\end{proof}
\paragraph{复习提示}上述表达式未必唯一，所以一般只能写上界。去掉交换性后，有限阶生成元之间交替相乘可能产生无限多个元素。

\subsection{习题10：两组对换生成对称群}
\paragraph{题目}证明$S_n$由$(12),(13),\ldots,(1n)$生成，也由$(12),(23),\ldots,(n-1\ n)$生成。
\paragraph{考点与思路}每个置换都是对换的乘积；因此只需把任意对换表示成指定生成元的乘积。
\textbf{原文依据：}现有笔记本题解答，以及正文“每个置换是对换的积”的命题。
\paragraph{所用结论}\toolref{a-permutation-generation}、\toolref{a-cycle-conjugation}。
\begin{proof}[详细证明]
先考虑第一组。涉及点1的对换已经在列表中；若$i,j\ne1$且$i\ne j$，利用共轭公式
\[
u(ab)u^{-1}=(u(a)\ u(b))
\]
取$u=(1i)$、$(ab)=(1j)$，得
\[
(ij)=(1i)(1j)(1i).
\]
故第一组生成全部对换，进而生成$S_n$。

再考虑相邻对换。固定$i<j$。若$j=i+1$，$(ij)$已经是相邻对换。若$j>i+1$，令
\[
u=(i\ i+1)(i+1\ i+2)\cdots(j-2\ j-1).
\]
从右至左作用时，$j-1$逐步被送到$i$，而$j$不受影响，因此$u(j-1)=i$、$u(j)=j$。共轭公式给出
\[
(ij)=u(j-1\ j)u^{-1}.
\]
右边的$u$、$u^{-1}$和中间对换都由相邻对换组成。因此第二组也生成全部对换，从而生成$S_n$。
\end{proof}
\paragraph{复习提示}必须说明“生成任意对换”为什么足够；共轭公式还将用于习题11、12。

\subsection{习题11：两个生成元生成对称群}
\paragraph{题目}当$n\ge2$时，证明$S_n=\langle(12),(12\cdots n)\rangle$。
\paragraph{考点与思路}用长轮换的幂把$(12)$共轭成所有相邻对换。
\textbf{原文依据：}现有笔记本题解答，以及习题10的相邻对换生成结论。
\paragraph{所用结论}\toolref{a-cycle-conjugation}、\toolref{a-permutation-generation}。
\begin{proof}[详细证明]
令$c=(12\cdots n)$，$K=\langle(12),c\rangle\le S_n$。对$1\le i<n$，有$c^{i-1}(1)=i$、$c^{i-1}(2)=i+1$，故
\[
c^{i-1}(12)c^{-(i-1)}=(i\ i+1)\in K.
\]
于是$K$包含习题10中生成$S_n$的全部相邻对换，从而$S_n\le K$。结合$K\le S_n$，即得$K=S_n$。$n=2$时两个指定生成元相同，仍生成$S_2$。
\end{proof}
\paragraph{复习提示}生成集可以有重复元素；关键是其生成子群是否包含一组已知生成元。

\subsection{习题12：两个生成元生成交错群}
\paragraph{题目}当$n>2$为奇数时，证明
\[
A_n=\langle(123),(12\cdots n)\rangle;
\]
当$n>2$为偶数时，证明
\[
A_n=\langle(123),(23\cdots n)\rangle.
\]
\textbf{原题订正：}按教材书页98（PDF第108页）核对，原笔记将两个生成元抄为偶数时$(12\cdots n)^2$、奇数时$(34\cdots n)$；本题已恢复教材中的生成元。下面仍沿用原笔记的“共轭构造3-轮换，再逐次附加新点”的证明思路。
\paragraph{考点与思路}先由共轭取得一列3-轮换，再说明这些轮换涉及的点彼此连接时，为什么能生成全部偶置换。
\textbf{原文依据：}题目按教材书页98核正；证明沿用现有笔记本题解答及其前置的关于一列3-轮换逐次加入新点的补充引理，以及正文“每个偶置换是3-轮换的积”的命题。教材该处只有题目，没有给出证明。
\paragraph{所用结论}\toolref{a-cycle-conjugation}、\toolref{a-triple-attach}、\toolref{a-permutation-generation}；其中附加点引理的证明在本题先行展开。
\begin{proof}[详细证明]
先证明一个加入新点的结论。设$U$至少有三个点，某子群$K$已包含在$U$上作用、其余点固定的全部偶置换。

若$K$还含$(abz)$，其中$a,b\in U$、$z\notin U$，则$K$包含任意$(uvz)$，其中$u,v\in U$且$u\ne v$。理由如下：当$|U|\ge4$时，先取把有序点对$(a,b)$送至$(u,v)$的置换；如果它是奇置换，在目标中除$u,v$以外的两个点上复合一个对换即可把它改成偶置换。对$(abz)$作共轭就得到$(uvz)$。当$|U|=3$时，$U$上的3-轮换循环移动点对，可得到三个有序点对；再对$(abz)$取逆，可交换前两个位置，得到其余三个有序点对。因此此时结论也成立。

任何$U\cup\{z\}$上的3-轮换，要么只涉及$U$，要么涉及$z$和$U$中的两个点。以上已经把两类全部包含，所以$K$包含$U\cup\{z\}$上的全部偶置换。

若新增的3-轮换只与$U$交一个点，写为$(cde)$，其中$c\in U$、$d,e\notin U$。按照现有笔记的附加点引理，选$u\in U\setminus\{c\}$，用$U$内把$c$送到$u$的偶置换共轭，得到$(ude)\in K$。这样的偶置换存在：在$U$中另取不同于$c,u$的一点，选适当方向的3-轮换即可。

现在先把$\{c,d,e\}$当成已有点集：$(cde)$的幂已经给出该三点集上的全部偶置换，而$(ude)$含有其中的$d,e$两个点。由前面加入一个新点的结论，$K$包含$\{c,d,e,u\}$上的全部偶置换，特别包含$(cud)$与$(cue)$。

回到原点集$U$，轮换$(cud)$中$c,u$是旧点，故先把$d$加入，得到$U\cup\{d\}$上的全部偶置换；再用$(cue)$把$e$加入，得到$U\cup\{d,e\}$上的全部偶置换。因此，\textbf{从一个3-轮换起，逐次加入与已有点集有交的3-轮换，就能生成最终点集上的全部偶置换}。

\textbf{奇数$n$的情形。}令$c=(12\cdots n)$、$K=\langle(123),c\rangle$。轮换$c$长为奇数，符号为$(-1)^{n-1}=1$，故两个生成元均为偶置换，$K\le A_n$。对$j=0,1,\ldots,n-3$，共轭公式给出
\[
c^j(123)c^{-j}=(j+1\ j+2\ j+3)\in K.
\]
这些轮换依次为$(123),(234),\ldots,(n-2\ n-1\ n)$。从三点集$\{1,2,3\}$开始，下一轮换的前两个点已经出现，第三个点是新点，所以反复使用前面的附加一个新点结论，即得到全部$n$个点上的偶置换。因此$A_n\le K$，结合反向包含，得$K=A_n$。当$n=3$时，$c=(123)$，它本身已经生成$A_3$。

\textbf{偶数$n$的情形。}此时$n\ge4$。令$d=(23\cdots n)$、$K=\langle(123),d\rangle$。轮换$d$长为奇数$n-1$，符号为$(-1)^{n-2}=1$，故$K\le A_n$。

$d$固定1。对$j=0,1,\ldots,n-3$，有$d^j(2)=j+2$、$d^j(3)=j+3$，因此
\[
d^j(123)d^{-j}=(1\ j+2\ j+3)\in K.
\]
这些轮换依次为$(123),(134),(145),\ldots,(1\ n-1\ n)$。从$\{1,2,3\}$起，每次轮换中的1和另一个点属于已有点集，第三个点是新点。因此仍逐次用附加一个新点的结论，得到全部$n$个点上的偶置换，即$A_n\le K$。所以$K=A_n$。
\end{proof}
\paragraph{复习提示}只说明生成元都是偶置换，只能得到生成子群包含于$A_n$。还须证明它足够大；这里的关键是将共轭构造与加入新点的论证结合。

\subsection{习题15：商群中子群的交}
\paragraph{题目}设$N\trianglelefteq G$、$H_1,H_2\le G$，且$N\le H_1\cap H_2$。证明
\[
(H_1/N)\cap(H_2/N)\cong(H_1\cap H_2)/N.
\]
原书要求同构；现有笔记通过陪集逐项检验证明两者在$G/N$中实际相等，以下保留此原有方法，从而得到所求同构。
\paragraph{考点与思路}在商群$G/N$中逐个陪集检验。包含$N$的条件保证一个代表元在子群内时，整个陪集都在其中。
\textbf{原文依据：}现有笔记本题解答（逐个陪集检验）。
\paragraph{所用结论}\toolref{a-quotient-membership}、\toolref{g-normal}。
\begin{proof}[详细证明]
因为$N\trianglelefteq G$且$N\le H_i$，$N$在每个$H_i$中也正规，各商群均有定义。对任意$g\in G$，证明
\[
gN\in H_i/N\iff g\in H_i.
\]
若$g\in H_i$，这是$H_i/N$的定义。反过来，若$gN=hN$，其中$h\in H_i$，则$h^{-1}g\in N\le H_i$，所以$g=h(h^{-1}g)\in H_i$。

因此对任意$gN\in G/N$，有
\begin{align*}
gN\in(H_1/N)\cap(H_2/N)
&\iff g\in H_1\text{ 且 }g\in H_2\\
&\iff g\in H_1\cap H_2\\
&\iff gN\in(H_1\cap H_2)/N.
\end{align*}
两边由同样的陪集组成，故相等。
\end{proof}
\paragraph{复习提示}结论是在$G/N$中的实际相等，比“同构”更强。不能遗漏$N\le H_i$，否则$H_i/N$甚至未必有定义。

\subsection{习题29：大子集的乘积}
\paragraph{题目}设$G$为有限群，$A,B\subseteq G$，且$|A|+|B|>|G|$。证明$AB=G$，其中$AB=\{ab:a\in A,b\in B\}$。
\paragraph{考点与思路}要把任意$g$写成$ab$，等价于让$A$和$gB^{-1}$有交点。
\textbf{原文依据：}现有笔记本题解答（集合交与计数）。
\paragraph{所用结论}\toolref{g-laws}（取逆和左乘均为双射）、\toolref{n-binomial-count}（有限集合交非空判据）。
\begin{proof}[详细证明]
固定$g\in G$。取逆映射和左乘$g$的映射都是双射，所以$|gB^{-1}|=|B|$。若$A\cap gB^{-1}=\varnothing$，它们的并集作为$G$的子集应满足
\[
|A|+|B|=|A\cup gB^{-1}|\le|G|,
\]
与题设矛盾。故存在$a\in A\cap gB^{-1}$。

由$a\in gB^{-1}$，可取$b\in B$使$a=gb^{-1}$，右乘$b$得$g=ab\in AB$。因为$g$任意，$G\subseteq AB$。又$A,B\subseteq G$且群运算封闭，$AB\subseteq G$，所以$AB=G$。
\end{proof}
\paragraph{复习提示}严格大于不能改成大于等于。例如在二阶群$\{e,a\}$中取$A=B=\{e\}$，有$|A|+|B|=|G|$，但$AB\ne G$。

\subsection{习题30：真子群的共轭不能覆盖有限群}
\paragraph{题目}设$G$为有限群，$H<G$为真子群。证明
\[
G\ne\bigcup_{g\in G}gHg^{-1}.
\]
\paragraph{考点与思路}共轭子群的数目至多为$[G:H]$；它们都包含单位元，因此并集的大小比直接相加至少少掉若干次重复。
\textbf{原文依据：}现有笔记本题解答，以及群作用中的轨道与稳定子计数。
\paragraph{所用结论}\toolref{a-orbit-stabilizer}、\toolref{g-cosets}（指数关系）、\toolref{n-binomial-count}（共享单位元的并集上界）。
\begin{proof}[详细证明]
记$|G|=n$、$|H|=h$、$m=[G:H]=n/h>1$。$G$通过共轭作用在$H$的共轭子群集合上。稳定$H$的元素正是$N_G(H)$，所以不同共轭子群的数目为
\[
r=[G:N_G(H)].
\]
对每个$h_0\in H$，有$h_0Hh_0^{-1}=H$，故$H\le N_G(H)$，由指数公式得到$r\le m$。

把不同共轭子群列为$H_1,\ldots,H_r$。每个$H_i$有$h$个元素且含$e$，因此
\begin{align*}
\left|\bigcup_{g\in G}gHg^{-1}\right|
&=1+\left|\bigcup_{i=1}^r(H_i\setminus\{e\})\right|\\
&\le1+\sum_{i=1}^r|H_i\setminus\{e\}|\\
&=1+r(h-1)\\
&\le1+m(h-1)=n-m+1<n.
\end{align*}
并集元素数严格小于$|G|$，故不可能覆盖$G$。当$H=\{e\}$时，该估计仍成立。
\end{proof}
\paragraph{复习提示}必须按“不同共轭子群”计数，不能把每个$g$都当成给出一个新子群。有限性用于基数估计。

\subsection{习题31：对称群四阶情形的西罗子群}
\paragraph{题目}求$S_4$的全部西罗2子群与西罗3子群。
\paragraph{考点与思路}8阶子群可由两两分块的稳定子构造；3阶子群由3-轮换生成，再用西罗数条件证明没有遗漏。
\textbf{原文依据：}现有笔记本题解答（分块稳定子与3-轮换计数），西罗第三定理。
\paragraph{所用结论}\toolref{a-orbit-stabilizer}（稳定子为子群）、\toolref{a-sylow-existence}、\toolref{a-sylow-count}、\toolref{g-lagrange}、\toolref{n-binomial-count}、\toolref{a-cycle-conjugation}（不相交轮换分解与置换的阶）。
\begin{proof}[详细证明]
因$|S_4|=24=2^3\cdot3$，西罗2子群的阶为8，西罗3子群的阶为3。

\textbf{西罗2子群。}将四个点分成两个无序二元组，有三种分块
\[
\mathcal P_1=\{\{1,2\},\{3,4\}\},\quad
\mathcal P_2=\{\{1,3\},\{2,4\}\},\quad
\mathcal P_3=\{\{1,4\},\{2,3\}\}.
\]
令$P_i=\{\sigma\in S_4:\sigma(\mathcal P_i)=\mathcal P_i\}$，即保持分块、允许交换两块的置换。稳定子在复合及取逆下封闭，故$P_i$是子群。可选择是否交换两块，再分别选择每块内部的对应，共有$2\cdot2\cdot2=8$种选择，所以$|P_i|=8$，均为西罗2子群。

三者彼此不同：$(12)\in P_1$而$(12)\notin P_2,P_3$；$(13)\in P_2$而$(13)\notin P_3$。又$n_2\mid3$，因此已有的三个子群正是全部。它们也可写为
\[
P_1=\langle(1324),(12)\rangle,\quad
P_2=\langle(1234),(13)\rangle,\quad
P_3=\langle(1243),(14)\rangle.
\]
以中间一组为例，$r=(1234)$、$s=(13)$满足$r^4=s^2=e$、$srs=r^{-1}$；八个$r^k,sr^k$互异，故它们生成整个$P_2$。其余两组按点的重新命名得到。因此三个子群都同构于正方形的8阶二面体群$D_4$。

\textbf{西罗3子群。}$S_4$的3阶元素只能是3-轮换。先选三个移动点，再选两种轮换方向，共有$\binom43\cdot2=8$个。一个3阶循环群恰有两个非单位元，且不同3阶子群交仅含$e$，因此恰有四个：
\[
\langle(123)\rangle,\quad\langle(124)\rangle,\quad
\langle(134)\rangle,\quad\langle(234)\rangle.
\]
\end{proof}
\paragraph{复习提示}西罗子群的“阶”与“个数”分开计算；本题分别是阶8、个数3，以及阶3、个数4。

\subsection{习题32：56阶与148阶群不单}
\paragraph{题目}证明56阶群和148阶群都不是单群。
\paragraph{考点与思路}148阶直接求出唯一37阶子群；56阶若7阶子群不唯一，就用元素计数迫使西罗2子群唯一。
\textbf{原文依据：}现有笔记本题解答，西罗第三定理。
\paragraph{所用结论}\toolref{a-sylow-existence}、\toolref{a-sylow-count}、\toolref{g-lagrange}、\toolref{n-binomial-count}。
\begin{proof}[详细证明]
\textbf{148阶。}设$|G|=148=4\cdot37$。西罗定理给出
\[
n_{37}\mid4,\qquad n_{37}\equiv1\pmod{37}.
\]
4的正因子只有1、2、4，只有1满足同余条件，所以$n_{37}=1$。唯一的37阶子群正规，且其阶既不为1也不为148，故$G$不单。

\textbf{56阶。}设$|G|=56=8\cdot7$。由$n_7\mid8$、$n_7\equiv1\pmod7$，得$n_7=1$或8。若$n_7=1$，唯一7阶子群就是非平凡真正规子群。

设$n_7=8$。不同7阶子群的交是各自的子群，阶整除7；若交非平凡，两者就相等。因此它们的非单位元互不重复，共有
\[
8(7-1)=48
\]
个7阶元素。全群除去这些元素后只剩8个元素，其中含$e$。

任取西罗2子群$P$，它有8个元素。由拉格朗日定理，$P$中元素阶只能为1、2、4、8，不能为7。因此$P$包含于上述剩余的8元集合。两者大小相同，必相等。每个西罗2子群都只能是该集合，所以$n_2=1$，$P$正规。又$1<8<56$，故$G$不单。
\end{proof}
\paragraph{复习提示}不能在知道$n_7=8$后就断言不单；真正的正规子群是进一步计数得到的唯一西罗2子群。

\subsection{习题36：一般线性群的西罗子群}
\paragraph{题目}设$p$为素数、$n\ge1$，求$\GL_n(\F_p)$的一个西罗$p$子群，并确定其正规化子及西罗$p$子群的个数。
\paragraph{考点与思路}先构造单位上三角群并计算其阶，再用有序基计数核对它恰好具有群阶中的全部$p$因子；通过公共不动向量逐层恢复标准子空间链，确定正规化子。
\textbf{原文依据：}现有笔记本题解答（单位上三角群、可逆矩阵计数、公共不动空间），西罗共轭定理。
\paragraph{所用结论}\toolref{g-basis-count}、\toolref{a-sylow-existence}、\toolref{a-sylow-conjugacy}、\toolref{a-sylow-count}、\toolref{a-orbit-stabilizer}；本题实际构造与正规化子计算汇总为\toolref{a-gl-sylow}，不是未经证明先套用此结论。
\begin{proof}[详细证明]
记$G=\GL_n(\F_p)$，按照现有笔记先取
\[
U=\{(a_{ij}):a_{ii}=1,\ a_{ij}=0\ (i>j)\}.
\]
两个这样的矩阵之积仍为上三角且对角线全1。若$u=I+T\in U$，则$T$严格上三角，故$T^n=0$，并且
\[
u^{-1}=I-T+T^2-\cdots+(-1)^{n-1}T^{n-1}\in U.
\]
因而$U$是子群。严格上三角位置有$n(n-1)/2$个，每个位置可独立取$p$个值，故$|U|=p^{n(n-1)/2}$。

接着展开现有笔记所用的可逆矩阵计数公式。可逆矩阵的列组成有序基：第一列有$p^n-1$种选择；已选$k$个线性无关列时，它们张成的空间有$p^k$个向量，下一列有$p^n-p^k$种选择。因此
\begin{align*}
|G|&=\prod_{k=0}^{n-1}(p^n-p^k)\\
&=p^{0+1+\cdots+(n-1)}\prod_{i=1}^n(p^i-1)\\
&=p^{n(n-1)/2}\prod_{i=1}^n(p^i-1).
\end{align*}
每个$p^i-1$均不被$p$整除，所以$|G|$中最大的$p$幂因子是$p^{n(n-1)/2}$，恰好等于$|U|$，故$U$为西罗$p$子群。

令$B$为全体可逆上三角矩阵组成的群。上三角矩阵之积及逆仍为上三角，对角线在乘法下逐项相乘。因此$bub^{-1}$的对角线全1，得到$B\le N_G(U)$。

反过来，设$g\in N_G(U)$。当$n\ge2$时，以$e_i$表示第$i$个标准基向量，$E_{ij}$表示仅第$(i,j)$分量为1的矩阵。对$i<j$有$I+E_{ij}\in U$，且
\[
(I+E_{ij})v=v+v_je_i.
\]
若$v$被$U$的每个元素固定，取$i=1$、$j=2,\ldots,n$，就得$v_2=\cdots=v_n=0$。反过来$e_1$的所有倍数确实被全部$U$固定。因此公共不动空间为$L_1=\langle e_1\rangle$。

正规化条件保证$g$保持这个空间：若$v\in L_1$、$u\in U$，则$g^{-1}ug\in U$，所以
\[
u(gv)=g(g^{-1}ug)v=gv.
\]
于是$g(L_1)=L_1$。在商空间$\F_p^n/L_1$上，$U$的像恰为关于基$e_2+L_1,\ldots,e_n+L_1$的单位上三角群，而$g$诱导的映射仍正规化这个像。重复上述公共不动空间论证，或对维数归纳，可得$g$保持整个链
\[
L_1\subset L_2\subset\cdots\subset L_n,
\qquad L_i=\langle e_1,\ldots,e_i\rangle.
\]
这意味着第$i$列$g(e_i)$只能出现在前$i$个坐标位置，因此$g$上三角，即$g\in B$。所以$N_G(U)=B$。

$B$的每个对角线分量有$p-1$种选择，严格上三角分量仍有$p^{n(n-1)/2}$种选择，故
\[
|B|=(p-1)^np^{n(n-1)/2}.
\]
由共轭作用和西罗共轭定理，全部西罗$p$子群数为
\[
n_p=[G:N_G(U)]
=\frac{\prod_{i=1}^n(p^i-1)}{(p-1)^n}.
\]
当$n=1$时$U=\{1\}$、$B=G$，正规化子与计数公式仍成立。
\end{proof}
\paragraph{复习提示}构造了具有正确阶的子群才算证明存在一个西罗子群。求个数时应除以正规化子阶，而非直接除以西罗子群阶。

\subsection{习题38：素数阶正规子群位于中心}
\paragraph{题目}设$G$为有限$p$群，$N\trianglelefteq G$且$|N|=p$，证明$N\le Z(G)$。
\paragraph{考点与思路}共轭给出从$p$群到$p-1$阶自同构群的同态，其像的阶只能为1。
\textbf{原文依据：}现有笔记本题解答，以及正文循环群自同构群的命题。
\paragraph{所用结论}\toolref{a-cycle-automorphisms}、\toolref{g-first-iso}、\toolref{g-lagrange}、\toolref{n-gcd}。
\begin{proof}[详细证明]
因$N$正规，每个$g\in G$都通过共轭给出$N$的自同构。定义
\[
\rho:G\longrightarrow\Aut(N),\qquad
\rho(g)(x)=gxg^{-1}\quad(x\in N).
\]
共轭映射可逆且保持乘法；又$c_gc_h=c_{gh}$，故$\rho$为同态。

$N$阶为素数，任一非单位元$a$均生成$N$。自同构由$a$的像唯一决定；像可以是$a,a^2,\ldots,a^{p-1}$中任意一个，每个选择都给出自同构。因此$|\Aut(N)|=p-1$。

由第一同构定理及拉格朗日定理，$|\rho(G)|$既整除$|G|=p^r$，又整除$p-1$。这两个整数互素，所以$|\rho(G)|=1$。故$gxg^{-1}=x$对所有$g\in G,x\in N$成立，也就是$xg=gx$。因此$N\le Z(G)$。
\end{proof}
\paragraph{复习提示}正规性本身只保证共轭后仍在$N$；这里还利用阶数互素，才把“仍在$N$”加强为“每个元素都不变”。

\subsection{习题39：中心商群循环推出交换}
\paragraph{题目}若$G/Z(G)$是循环群，证明$G$为交换群。
\paragraph{考点与思路}同一个陪集生成元的幂提供非中心部分，其余因子都可自由交换。
\textbf{原文依据：}现有笔记本题解答（中心陪集表示与交换计算）。
\paragraph{所用结论}\toolref{g-cosets}与\toolref{g-laws}；本题证明的结论汇总为\toolref{a-center-quotient}。
\begin{proof}[详细证明]
设$G/Z(G)=\langle aZ(G)\rangle$。任意$x\in G$的陪集都有形式$xZ(G)=a^mZ(G)$，其中$m\in\Z$，故$a^{-m}x\in Z(G)$。因此任意$x,y\in G$可写成
\[
x=a^mz,\qquad y=a^nw,
\qquad z,w\in Z(G),\ m,n\in\Z.
\]
中心元与所有元素交换，故
\begin{align*}
xy&=a^mza^nw\\
&=a^ma^nzw\\
&=a^{m+n}zw\\
&=a^{n+m}wz\\
&=a^nwa^mz=yx.
\end{align*}
$x,y$任意，所以$G$交换。上述论证不要求$G$有限，也包括商群平凡的情况。
\end{proof}
\paragraph{复习提示}结论实际上进一步给出$Z(G)=G$，所以中心商群若循环，就只能是平凡群。

\subsection{习题40：有限群嵌入交错群}
\paragraph{题目}证明任意有限群都同构于某个交错群的子群。
\paragraph{考点与思路}先用左乘作用嵌入对称群，再添两个点，用一个额外对换补偿奇偶性。
\textbf{原文依据：}现有笔记本题解答（凯莱嵌入后添两点补偿奇偶性）。
\paragraph{所用结论}\toolref{a-cayley}、\toolref{a-cycle-conjugation}（符号同态及不相交支集交换）、\toolref{g-first-iso}的核平凡单射判据。
\begin{proof}[详细证明]
设$|G|=m$。对每个$g\in G$，定义左乘置换$L_g:G\to G$、$x\mapsto gx$。$L_{g^{-1}}$为其逆，且$L_gL_h=L_{gh}$，所以$g\mapsto L_g$是到$m$点对称群的同态。若$L_g$为恒等映射，则$g=L_g(e)=e$，因此它单射。给$G$的元素编号后得到$G\hookrightarrow S_m$。

把$S_m$中置换延拓到$m+2$个点，令其固定$m+1,m+2$。令$t=(m+1\ m+2)$。对$\sigma\in S_m$，设$\varepsilon(\sigma)\in\{0,1\}$表示奇偶性，偶时为0，奇时为1，定义
\[
\iota(\sigma)=\sigma t^{\varepsilon(\sigma)}.
\]
因为$t$是奇置换，$\iota(\sigma)$总是偶置换，即属于$A_{m+2}$。

$t$只作用于最后两个点，因而与$S_m$中的元素交换。又奇偶性在复合下模2相加，所以
\begin{align*}
\iota(\sigma)\iota(\tau)
&=\sigma t^{\varepsilon(\sigma)}\tau t^{\varepsilon(\tau)}\\
&=\sigma\tau t^{\varepsilon(\sigma)+\varepsilon(\tau)}\\
&=\sigma\tau t^{\varepsilon(\sigma\tau)}
=\iota(\sigma\tau).
\end{align*}
若$\iota(\sigma)=e$，将其限制于前$m$个点，就有$\sigma=e$，故$\iota$单射。复合$G\hookrightarrow S_m\hookrightarrow A_{m+2}$即得所求。
\end{proof}
\paragraph{复习提示}不能直接说$S_m\le A_{m+2}$；必须指定这个保持乘法的嵌入。额外对换恰好抵消原置换的奇偶性。

\subsection{习题41：与导群交平凡的正规子群}
\paragraph{题目}若$N\trianglelefteq G$且$N\cap G^{(1)}=\{e\}$，证明$N\le Z(G)$。
\paragraph{考点与思路}任取一个$N$中元素与任意群元的换位子，证明它同时落在$N$和导群中。
\textbf{原文依据：}现有笔记本题解答（换位子同时属于正规子群与导群）。
\paragraph{所用结论}\toolref{g-normal}、\toolref{g-subgroup}、\toolref{a-derived-functor}中的换位子群定义。
\begin{proof}[详细证明]
取$n\in N$、$g\in G$。正规性给出$g^{-1}ng\in N$；因为$n^{-1}\in N$且$N$对乘法封闭，
\[
[n,g]=n^{-1}g^{-1}ng=n^{-1}(g^{-1}ng)\in N.
\]
另一方面，$G^{(1)}$由全部换位子生成，故$[n,g]\in G^{(1)}$。于是$[n,g]\in N\cap G^{(1)}=\{e\}$，即
\[
n^{-1}g^{-1}ng=e.
\]
左乘$n$再左乘$g$得到$ng=gn$。因为$g$任意，$n\in Z(G)$；因为$n$任意，$N\le Z(G)$。
\end{proof}
\paragraph{复习提示}正规性用于证明换位子属于$N$，交平凡用于证明换位子等于$e$，两个条件分别承担不同步骤。

\subsection{习题42：可解群的子群与商群}
\paragraph{题目}证明可解群的子群与商群仍然可解。
\paragraph{考点与思路}子群的导出列包含在原群的导出列中；满同态则把导出列逐项映满到像群的导出列。
\textbf{原文依据：}现有笔记本题解答，以及正文“可解群的子群、同态像及扩张”命题的前两部分。
\paragraph{所用结论}\toolref{a-derived-functor}；本题据此证明的可解性结论汇总为\toolref{a-solvable-closure}。
\begin{proof}[详细证明]
设$G^{(r)}=\{e\}$。

\textbf{子群。}若$H\le G$，则$H$中任意两个元素的换位子也是$G$中的换位子，因此$H^{(1)}\le G^{(1)}$。逐次使用同样的理由，归纳得到$H^{(i)}\le G^{(i)}$对所有$i\ge0$成立。特别地，
\[
H^{(r)}\le G^{(r)}=\{e\},
\]
所以$H$可解。

\textbf{商群。}设$N\trianglelefteq G$，令$\pi:G\twoheadrightarrow G/N$为自然满同态。任意同态都满足
\[
\pi([x,y])=[\pi(x),\pi(y)].
\]
这给出$\pi(G^{(1)})\le(G/N)^{(1)}$。反过来，商群任意两个元素都能取原像，故商群的每个换位子都是原群一个换位子的像；这些换位子生成商群导群，所以反向包含也成立。于是$\pi(G^{(1)})=(G/N)^{(1)}$。

把此结论反复应用于$\pi$在各导群上的限制，得到
\[
(G/N)^{(i)}=\pi(G^{(i)}).
\]
取$i=r$，右边为$\{N\}$，故$G/N$可解。
\end{proof}
\paragraph{复习提示}子群情形一般是包含关系；商群情形因映射满而得到相等关系。不能把“导群元素”都误认为单个换位子，它们还包括换位子的有限乘积。

\subsection{习题43：不同素数的平方乘积阶群可解}
\paragraph{题目}设$p,q$为不同素数，$|G|=p^2q$，证明$G$可解。
\paragraph{考点与思路}沿用现有笔记的顺序：由习题25得到正规西罗子群，由习题20判断它与商群均交换，最后应用可解群的扩张性质。下面将所依赖的两道习题及扩张性质一并展开。
\textbf{原文依据：}现有笔记本题解答、习题25与20，以及正文“可解群的子群、同态像及扩张”命题的扩张部分。
\paragraph{所用结论}\toolref{a-p-square-q-normal}、\toolref{a-p-square-abelian}、\toolref{a-solvable-extension}；前两项所需的\toolref{a-sylow-count}、\toolref{a-class-equation}、\toolref{a-center-quotient}也在本题展开。
\begin{proof}[详细证明]
\textbf{第一步：至少有一个西罗子群正规。}设$n_p,n_q$分别为西罗子群数。若$n_p=1$或$n_q=1$，已经得到结论。反设两者都大于1。由$n_p\mid q$，必有$n_p=q$；由$n_p\equiv1\pmod p$，得$q>p$。

又$n_q\mid p^2$且$n_q>1$，所以$n_q=p$或$p^2$。因为$p<q$，整数$p$不可能模$q$余1，故$n_q=p^2$。于是
\[
q\mid p^2-1=(p-1)(p+1).
\]
$q$为素数，故整除其中一个因子。$q>p$排除$q\mid p-1$，所以$q\mid p+1$。由于$p<q\le p+1$，有$q=p+1$。相邻两个整数都是素数，只可能$p=2,q=3$。

因此只剩$|G|=12$且$n_2=3,n_3=4$的情形。四个3阶子群互交仅含$e$，共含$4\cdot2=8$个3阶元素。其余元素只有四个，含$e$。每个西罗2子群有四个元素，且不能含3阶元素，所以必等于这四元集合，因而西罗2子群唯一，与$n_2=3$矛盾。第一步得证。

\textbf{第二步：阶为$p^2$的群交换。}先回顾$p$群的类方程：非中心共轭类的大小是大于1的$p$幂，因此
\[
|P|=|Z(P)|+\sum\text{非中心共轭类大小}
\]
说明非平凡有限$p$群的中心阶被$p$整除，中心非平凡。若$|P|=p^2$，则$|Z(P)|=p$或$p^2$。若为$p$，商群$P/Z(P)$阶为$p$，故循环；由习题39，$P$交换，反而有$Z(P)=P$，与中心阶为$p$矛盾。所以必有$Z(P)=P$，$P$交换。

\textbf{第三步：应用可解群的扩张性质。}取第一步给出的正规西罗子群$N$。若$|N|=p^2$，第二步说明$N$交换，而$G/N$为$q$阶循环群；若$|N|=q$，则$N$循环，而$G/N$阶为$p^2$，同样交换。因此两种情况下$N$与$G/N$均可解。

现有笔记“可解群的子群、同态像及扩张”命题说明：正规子群与相应商群均可解时，原群可解。在这里使用该命题，即得$G$可解。为补全使用理由，按该命题的原证明，设$N^{(s)}=\{e\}$、$(G/N)^{(r)}=\{N\}$。对自然满同态$\pi:G\to G/N$，习题42的导群像公式给出
\[
\pi(G^{(r)})=(G/N)^{(r)}=\{N\},
\]
所以$G^{(r)}\le N$。再逐次取$s$次导群，得到
\[
G^{(r+s)}=\bigl(G^{(r)}\bigr)^{(s)}\le N^{(s)}=\{e\}.
\]
这正是扩张性质推出$G$可解的过程。本题中$N$和$G/N$均交换，可取$r=s=1$。
\end{proof}
\paragraph{复习提示}本题并未要求$p<q$；证明里只有在反设两个西罗子群都不正规时，才推出$q>p$。应特别记住12阶的计数补充。

\subsection{习题44：有限素数幂阶群可解}
\paragraph{题目}证明任意有限$p$群都是可解群。
\paragraph{考点与思路}沿用现有笔记的中心商群归纳：类方程给出非平凡中心，商去中心后阶数下降；中心与中心商群均可解，再应用可解群的扩张性质。
\textbf{原文依据：}现有笔记本题解答、类方程，以及正文“可解群的子群、同态像及扩张”命题的扩张部分。
\paragraph{所用结论}\toolref{n-induction}、\toolref{a-class-equation}、\toolref{a-solvable-extension}。
\begin{proof}[详细证明]
对群阶归纳。平凡群交换，可解。设结论对阶小于$|G|$的$p$群成立，且$G$非平凡。

由习题43中类方程的论证，$Z(G)\ne\{e\}$。若$Z(G)=G$，则$G$交换，$G^{(1)}=\{e\}$，结论成立。

以下设$Z(G)$为真子群。因$|Z(G)|>1$，商群$G/Z(G)$仍是$p$群，且阶严格小于$|G|$。由归纳假设，$G/Z(G)$可解。另一方面，$Z(G)$交换，故也可解；中心又是$G$的正规子群。因此对$N=Z(G)$应用可解群的扩张性质，就得到$G$可解，归纳完成。

这里的扩张性质正是习题43第三步按现有笔记命题所展开的结论：若$(G/Z(G))^{(r)}$平凡，则$G^{(r)}\le Z(G)$；由于$Z(G)^{(1)}=\{e\}$，再取一次导群即有$G^{(r+1)}=\{e\}$。这补明了原文“由正规子群与商群可解推出原群可解”的理由。
\end{proof}
\paragraph{复习提示}$p$群不一定交换，但中心一定非平凡。归纳所需的是商群阶数变小，不能仅说“商群仍是$p$群”就直接套结论。

\subsection{习题45：两个可解商群与交核}
\paragraph{题目}若$H,K\trianglelefteq G$且$G/H,G/K$均可解，证明$G/(H\cap K)$可解。
\paragraph{考点与思路}将两种商映射合并成到直积的映射，其核正好是交。再证明两个可解群的直积可解。
\textbf{原文依据：}现有笔记本题解答（到两个商群直积的同态、核与导出列），第一同构定理。
\paragraph{所用结论}\toolref{g-first-iso}、\toolref{a-solvable-product}、\toolref{a-solvable-closure}；直积导群公式按\toolref{a-derived-functor}的定义在本题计算。
\begin{proof}[详细证明]
定义
\[
\varphi:G\longrightarrow(G/H)\times(G/K),
\qquad g\longmapsto(gH,gK).
\]
两坐标都是同态，所以$\varphi$为同态。且
\[
g\in\ker\varphi\iff gH=H\text{ 且 }gK=K
\iff g\in H\cap K.
\]
第一同构定理给出
\[
G/(H\cap K)\cong\varphi(G)\le(G/H)\times(G/K).
\]

为证明右边的直积可解，记$A=G/H$、$B=G/K$。逐分量计算有
\[
[(a,b),(a',b')]=([a,a'],[b,b']).
\]
故$(A\times B)^{(1)}\le A^{(1)}\times B^{(1)}$。反向，取第二坐标为$e$可得到$([a,a'],e)$；取第一坐标为$e$可得到$(e,[b,b'])$，这些元素的乘积生成$A^{(1)}\times B^{(1)}$。因此两边相等。归纳得
\[
(A\times B)^{(i)}=A^{(i)}\times B^{(i)}.
\]
若$A^{(r)}=\{e\}$、$B^{(s)}=\{e\}$，取$t=\max\{r,s\}$，得到$(A\times B)^{(t)}=\{(e,e)\}$，直积可解。由习题42，可解群的子群$\varphi(G)$可解，同构将此性质传给$G/(H\cap K)$。
\end{proof}
\paragraph{复习提示}$\varphi$未必满到整个直积，因此只写$G/(H\cap K)\cong\varphi(G)$，不能直接写成与整个直积同构。

\subsection{习题46：内自同构恰有两个的条件不可能发生}
\paragraph{题目}若群$G$只有两个内自同构，证明$G$交换。进一步说明这一前提实际上不可能发生。
\paragraph{考点与思路}利用$\operatorname{Inn}(G)\cong G/Z(G)$，将自同构数目转化为中心商群的阶数。
\textbf{原文依据：}现有笔记本题解答、内自同构群与中心商群同构的命题，以及习题39。
\paragraph{所用结论}\toolref{a-inner-automorphisms}、\toolref{g-cyclic}、\toolref{g-lagrange}（素数阶群循环）、\toolref{a-center-quotient}。
\begin{proof}[详细证明]
假设$|\operatorname{Inn}(G)|=2$。由本章开始说明的同构，$G/Z(G)$是2阶群；2阶群必循环，故习题39推出$G$交换。

然而，交换群中任意$g,x$都有$gxg^{-1}=x$，所以所有内自同构都是恒等映射，必有$|\operatorname{Inn}(G)|=1$。这与假设等于2矛盾。因此不存在内自同构恰有两个的群，而题目的条件命题成立。
\end{proof}
\paragraph{复习提示}“推出交换”后还可以继续检查前提：这里不是找到一类特殊交换群，而是证明前提本身不可能满足。

\subsection{习题47：至少存在两个自同构}
\paragraph{题目}设$G$为有限群，$|G|>2$。证明$|\Aut(G)|\ge2$；2阶群是非平凡有限群中的例外。
\paragraph{考点与思路}按是否交换、是否有阶大于2的元素分类，分别构造非恒等内自同构、取逆映射、交换基向量的线性自同构。
\textbf{原文依据：}现有笔记本题解答的三种分类构造。
\paragraph{所用结论}\toolref{a-inner-automorphisms}、\toolref{g-laws}、\toolref{a-binary-basis-swap}。
\begin{proof}[详细证明]
恒等映射始终是自同构，因此只需再构造一个非恒等自同构。

\textbf{第一种情况：$G$不交换。}存在$a\notin Z(G)$，因而有$x\in G$使$ax\ne xa$。内自同构$c_a:y\mapsto aya^{-1}$满足$c_a(x)\ne x$，所以不是恒等映射。

\textbf{第二种情况：$G$交换，且有元素$a$的阶大于2。}取逆映射$f(x)=x^{-1}$满足
\[
f(xy)=y^{-1}x^{-1}=x^{-1}y^{-1}=f(x)f(y).
\]
又$f^2$为恒等映射，故$f$双射，是自同构。因为$a^2\ne e$，有$a^{-1}\ne a$，所以$f$非恒等。

\textbf{第三种情况：$G$交换，且所有元素平方均为$e$。}将群运算记成加法，单位元记成0，则$x+x=0$对所有$x$成立。规定$\F_2$上的标量乘法为$0x=0$、$1x=x$。向量空间的分配律中唯一需要特别检查的是
\[
(1+1)x=0=x+x;
\]
其余标量组合均直接满足，因此$G$成为$\F_2$上的向量空间。

因$G$有限，可从非零向量逐次加入不在已有张成空间中的向量，有限步后得到一组基$v_1,\ldots,v_d$。每个向量唯一写成$\sum_i\varepsilon_iv_i$，$\varepsilon_i\in\{0,1\}$，故$|G|=2^d$。由$|G|>2$得$d\ge2$。定义
\[
f\left(\varepsilon_1v_1+\varepsilon_2v_2+\sum_{i=3}^d\varepsilon_iv_i\right)
=\varepsilon_1v_2+\varepsilon_2v_1+\sum_{i=3}^d\varepsilon_iv_i.
\]
基表示唯一，故$f$定义良好；它保持加法，且$f^2$为恒等映射，因此是群自同构。又$f(v_1)=v_2\ne v_1$，所以非恒等。

三种情况覆盖所有群，均已构造第二个自同构。若$|G|=2$，唯一非单位元必须被自同构固定，单位元也被固定，因此只有恒等自同构，确为例外。
\end{proof}
\paragraph{复习提示}取逆在一般非交换群中不是同态；先分出交换情形才能使用。所有元素平方为单位元时，取逆映射又恰是恒等，需要另作构造。

\subsection{习题49：两个小对称群是完全群}
\paragraph{题目}证明$S_3$、$S_4$都是完全群，即中心平凡且全部自同构均为内自同构。
\paragraph{考点与思路}先证明中心平凡以数出内自同构，再让任意自同构作用在一个小的、由群自身性质确定的集合上，给自同构总数一个同样的上界。
\textbf{原文依据：}现有笔记本题解答、习题6的中心计算、习题31的四个西罗3子群，以及正文完全群的定义。
\paragraph{所用结论}\toolref{a-inner-automorphisms}、\toolref{a-automorphism-invariants}、\toolref{a-cycle-conjugation}、\toolref{a-permutation-generation}；四个西罗3子群已在本册习题31求出。
\begin{proof}[详细证明]
\textbf{第一步：中心平凡。}事实上$n\ge3$时$Z(S_n)=\{e\}$。若$z\in Z(S_n)$，则对每个$i\ne j$，
\[
(z(i)\ z(j))=z(ij)z^{-1}=(ij),
\]
所以$z$保持每个二元子集$\{i,j\}$。固定$i$，取互不相同的$i,j,k$，则$z(i)$同时属于$\{i,j\}$与$\{i,k\}$，只能等于$i$。故$z=e$。

因此$\operatorname{Inn}(S_3)\cong S_3$有6个元素，$\operatorname{Inn}(S_4)\cong S_4$有24个元素。只需证明各自的全部自同构数不超过这两个数。

\textbf{第二步：$S_3$。}自同构保持元素阶，所以置换$S_3$中的三个2阶元素$(12),(13),(23)$。这给出同态
\[
\Aut(S_3)\longrightarrow\operatorname{Sym}\{(12),(13),(23)\}\cong S_3.
\]
若自同构固定这三个元素，因它们生成$S_3$，该自同构对任意生成元乘积也固定，故为恒等映射。因此此同态单射，$|\Aut(S_3)|\le6$。内自同构已经有6个，故全部自同构都是内自同构，$S_3$完全。

\textbf{第三步：$S_4$。}由习题31，$S_4$有四个西罗3子群。记$Q_i$为固定点$i$、在其余三个点上作3-轮换的那个3阶子群。令
\[
H_i=\{\sigma\in S_4:\sigma(i)=i\}.
\]
我们先证明$N_{S_4}(Q_i)=H_i$。若$g$正规化$Q_i$，它把$Q_i$中非单位3-轮换共轭为同一子群中的非单位3-轮换。共轭将移动点集合映为移动点集合，因此$g$保持$\{1,2,3,4\}\setminus\{i\}$，从而$g(i)=i$。反过来，若$g$固定$i$，它共轭一个移动另外三个点的3-轮换后，仍是这三个点上的3-轮换，只能为原轮换或其逆，所以正规化$Q_i$。等式得证。

自同构把西罗3子群互相置换，并满足
\[
\alpha(N_{S_4}(Q))=N_{S_4}(\alpha(Q)).
\]
故它们也把四个$H_i$互相置换，给出同态
\[
\Psi:\Aut(S_4)\longrightarrow\operatorname{Sym}\{H_1,H_2,H_3,H_4\}\cong S_4.
\]
四个$H_i$彼此不同：若$i\ne j$，取$k\ne i,j$，则$(jk)$固定$i$而不固定$j$，所以$H_i\ne H_j$。

下面证明$\Psi$单射。设$\alpha\in\ker\Psi$，即$\alpha(H_i)=H_i$对每个$i$成立。任取$g\in S_4$，共轭点稳定子得到
\[
gH_ig^{-1}=H_{g(i)}.
\]
对该式应用$\alpha$，并利用它固定各个$H_i$，有
\begin{align*}
H_{\alpha(g)(i)}
&=\alpha(g)H_i\alpha(g)^{-1}\\
&=\alpha(gH_ig^{-1})\\
&=\alpha(H_{g(i)})=H_{g(i)}.
\end{align*}
因四个$H_i$彼此不同，必有$\alpha(g)(i)=g(i)$对所有$i$成立。所以$\alpha(g)$与$g$是相同置换。$g$任意，故$\alpha$为恒等自同构，$\ker\Psi$平凡。

于是$|\Aut(S_4)|\le|S_4|=24$。结合已有24个内自同构，得到$\Aut(S_4)=\operatorname{Inn}(S_4)$。再结合第一步的中心平凡，$S_4$也是完全群。
\end{proof}
\paragraph{复习提示}证明自同构作用的核平凡是本题关键。仅说“自同构置换四个西罗3子群”，还不能直接断言自同构至多24个；必须排除作用不忠实的可能。

\clearpage% 题目依据：原笔记 chapters/Part 2/chapter3-exercises.tex。
\section{第三章：环论}

\subsection*{本章使用的定义与约定}
\begin{itemize}
  \item \textbf{幺环}是有乘法单位元 $1$ 的环；乘法不预设交换。若 $u,v\in R$ 满足 $uv=1$，称 $v$ 为 $u$ 的\textbf{右逆}，称 $u$ 为 $v$ 的左逆。若还满足 $vu=1$，则 $u$ 是\textbf{单位}，$v=u^{-1}$。
  \item 非零元素 $a$ 若存在非零 $b$ 使 $ab=0$，本资料称 $a$ 为\textbf{左零因子}；若存在非零 $b$ 使 $ba=0$，称 $a$ 为\textbf{右零因子}。两种情况至少发生一种，就称 $a$ 为\textbf{零因子}。不同教材可能对左右名称采用相反约定，考试时写出所用等式即可消除歧义。
  \item 若存在正整数 $m$ 使 $a^m=0$，称 $a$ 为\textbf{幂零元}；若 $e^2=e$，称 $e$ 为\textbf{幂等元}。零元同时满足这两种性质。
  \item \textbf{保幺环同态} $\eta:R\to R'$ 满足
  \[
  \eta(a+b)=\eta(a)+\eta(b),\qquad
  \eta(ab)=\eta(a)\eta(b),\qquad \eta(1)=1'.
  \]
  它还满足 $\eta(0)=0'$、$\eta(-a)=-\eta(a)$。\textbf{整环}指 $1\ne0$、乘法交换且没有零因子的幺环。
\end{itemize}
本章四题均沿用现有笔记本题解答的构造、论证顺序及反例，只补充计算、条件检查和中间理由。已核对第三章正文；这四题原解并未援引正文中某个具名同态定理，故不另换用定理证明，也不虚列定理编号。

\subsection{习题1：由\texorpdfstring{$1-ab$}{1-ab}可逆推出\texorpdfstring{$1-ba$}{1-ba}可逆}\label{exercise:3:1}
\paragraph{题目}
设幺环 $R$ 中 $a,b\in R$，且 $1-ab$ 可逆。证明 $1-ba$ 可逆。

\paragraph{考点与思路}
在非交换环中，不能将 $ab$ 与 $ba$ 直接互换。先记 $c=(1-ab)^{-1}$，再尝试把原来的逆元嵌入 $bca$，构造 $1-ba$ 的逆元。必须分别验证左右两个乘积。
\textbf{原文依据：}现有笔记第三章习题1解答：令 $c=(1-ab)^{-1}$，先验算 $(1-ba)(1+bca)$，再验算反向乘积。

\paragraph{所用结论}\toolref{r-ring-laws}、\toolref{r-unit-laws}。
\begin{proof}[详细证明]
由 $c=(1-ab)^{-1}$ 的定义，
\[
(1-ab)c=1,\qquad c(1-ab)=1.
\]
把两式分别展开，得到
\[
c-abc=1,\qquad c-cab=1.
\]
令 $d=1+bca$。先计算 $(1-ba)d$，使用结合律把含 $b$ 和 $a$ 的项提取出来：
\begin{align*}
(1-ba)(1+bca)
  &=1+bca-ba-babca\\
  &=1+b(c-1-abc)a\\
  &=1.
\end{align*}
最后一步用了 $c-abc=1$。

另一方向需要使用另一条逆元等式：
\begin{align*}
(1+bca)(1-ba)
  &=1-ba+bca-bcaba\\
  &=1+b(c-1-cab)a\\
  &=1,
\end{align*}
其中最后一步用了 $c-cab=1$。因此 $d$ 是 $1-ba$ 的双侧逆元，且
\[
\boxed{(1-ba)^{-1}=1+b(1-ab)^{-1}a.}
\]
\end{proof}

\paragraph{复习提示}
记住逆元的构造次序是 $1+bca$。证明中没有交换任意两个因子，因此适用于非交换幺环。将 $a,b$ 对换也得到逆命题，故 $1-ab$ 可逆当且仅当 $1-ba$ 可逆。

\subsection{习题2：右逆不唯一、零因子与非单位的等价性}\label{exercise:3:2}
\paragraph{题目}
设幺环 $R$ 中的元素 $u$ 有右逆。证明下列三个条件等价：
\begin{enumerate}
  \item $u$ 有多于一个右逆；
  \item $u$ 是零因子；
  \item $u$ 不是单位。
\end{enumerate}

\paragraph{考点与思路}
先固定一个右逆 $v$，即 $uv=1$。两个右逆相减后落入方程 $ut=0$ 的解集中；反过来，给 $v$ 加上这样的 $t$ 就能产生新的右逆。条件 $uv=1$ 还排除了非零元素在 $u$ 左侧乘出零的可能。
\textbf{原文依据：}现有笔记第三章习题2解答：先以右逆之差证明零因子条件的等价性，再以 $v+(vu-1)$ 为右逆证明唯一右逆推出单位。

\paragraph{所用结论}\toolref{r-ring-laws}、\toolref{r-unit-laws}。
\begin{proof}[详细证明]
固定 $v\in R$ 满足 $uv=1$。若 $R$ 是只有一个元素的零环，则 $u$ 是单位，只有一个右逆且不是零因子，三个条件均不成立，等价性仍成立。以下设 $1\ne0$，于是 $u\ne0$。

\textbf{第一步：证明条件 (1) 与条件 (2) 等价。}
若另有 $w\ne v$ 满足 $uw=1$，则 $w-v\ne0$，且
\[
u(w-v)=uw-uv=1-1=0.
\]
因此 $u$ 是左零因子，特别是零因子。

反过来，先注意到：对于任意 $x\in R$，若 $xu=0$，则
\[
x=x\cdot1=x(uv)=(xu)v=0.
\]
所以不存在非零 $x$ 使 $xu=0$。若 $u$ 是零因子，必有非零 $t$ 满足 $ut=0$。于是
\[
u(v+t)=uv+ut=1,
\]
而 $v+t\ne v$，故 $u$ 至少有两个右逆。

\textbf{第二步：证明条件 (1) 与条件 (3) 等价。}
若 $u$ 是单位，任取右逆 $w$，由 $uw=1$ 得
\[
w=u^{-1}(uw)=u^{-1}.
\]
因而右逆唯一。所以右逆不唯一必推出 $u$ 不是单位。

另一方面，设 $u$ 的右逆唯一。注意
\[
u(vu-1)=(uv)u-u=u-u=0.
\]
因此 $v+(vu-1)$ 也是右逆，因为
\[
u\bigl(v+vu-1\bigr)=uv+u(vu-1)=1.
\]
由右逆的唯一性，$v+vu-1=v$，在加法群中消去 $v$，得到 $vu=1$。结合 $uv=1$，可知 $u$ 是单位。

取这条结论的逆否命题，$u$ 不是单位便推出其右逆不唯一。至此三个条件两两等价。
\end{proof}

\paragraph{复习提示}
这里的前提“已经有一个右逆”不可遗漏。一般幺环中“不是单位”未必是零因子，例如整数环中的 $2$。本题还证明：有右逆的非单位只能以 $ut=0$ 的形式成为零因子，不能有非零 $x$ 满足 $xu=0$。

\subsection{习题3：右逆不唯一便有无穷多个右逆}\label{exercise:3:3}
\paragraph{题目}
设幺环中 $u$ 有多于一个右逆，证明 $u$ 有无穷多个右逆。

\paragraph{考点与思路}
由 $uv=1$ 但 $vu\ne1$ 构造非零元素 $e=1-vu$。它满足 $ue=0$、$ev=0$，适合用来给已有的右逆 $v$ 添加不同的修正项。不能仅写 $v+ne$，因为正特征下这些元素可能重复。
\textbf{原文依据：}现有笔记第三章习题3解答及其引用的习题2：构造 $e=1-vu$ 和 $v+eu^k$，再在假定两项相等后右乘 $v^j$ 导出矛盾。

\paragraph{所用结论}\toolref{r-ring-laws}、\toolref{r-unit-laws}、\toolref{n-induction}。并使用本章习题2（\ref{exercise:3:2}）已经证明的“右逆不唯一推出非单位”；不将本题结论作为前提。
\begin{proof}[详细证明]
取一个右逆 $v$，使 $uv=1$。由习题2，$u$ 不是单位。若 $vu=1$，则 $v$ 同时是左逆和右逆，$u$ 便是单位，矛盾。因此
\[
e=1-vu\ne0.
\]
由 $uv=1$，有
\begin{align*}
ue&=u(1-vu)=u-(uv)u=0,\\
ev&=(1-vu)v=v-v(uv)=0.
\end{align*}

对每个整数 $k\ge0$，令
\[
x_k=v+eu^k,
\]
这里 $u^0=1$。因为 $ue=0$，
\[
ux_k=uv+(ue)u^k=1,
\]
所以每个 $x_k$ 都是 $u$ 的右逆。

下面证明它们两两不同。先由 $uv=1$ 归纳得到，对每个整数 $m\ge0$，
\[
u^mv^m=1.
\]
归纳步骤为
\[
u^{m+1}v^{m+1}=u^m(uv)v^m=u^mv^m=1.
\]
于是当 $0\le i<j$ 时，
\[
u^iv^j=(u^iv^i)v^{j-i}=v^{j-i}.
\]

假设 $x_i=x_j$，其中 $i<j$。消去 $v$ 后，$eu^i=eu^j$。两边右乘 $v^j$，得到
\[
ev^{j-i}=e.
\]
但 $j-i\ge1$ 且 $ev=0$，所以
\[
ev^{j-i}=(ev)v^{j-i-1}=0,
\]
从而 $e=0$，与构造矛盾。因此 $x_0,x_1,x_2,\ldots$ 两两不同，$u$ 有无穷多个右逆。
\end{proof}

\paragraph{复习提示}
证明的核心是“构造右逆”和“证明构造出的元素互异”两个步骤。等式 $u^mv^m=1$ 来自逐对消去中间的 $uv$，绝不意味着 $u,v$ 可以交换。本题也说明有限幺环中，有右逆的元素必为单位。

\subsection{习题8：满同态下的元素性质与整环性质}\label{exercise:3:8}
\paragraph{题目}
设 $\eta:R\twoheadrightarrow R'$ 为保幺满同态。判断下列命题是否成立，并证明或举反例：
\begin{enumerate}
  \item[(i)] 幂零元、幂等元的像仍是相应元素；
  \item[(ii)] 零因子的像仍是零因子；
  \item[(iii)] 若 $R$ 为整环，则 $R'$ 为整环；
  \item[(iv)] 若 $R'$ 为整环，则 $R$ 为整环；
  \item[(v)] 单位的像仍是单位；
  \item[(vi)] 像为单位的元素必为单位。
\end{enumerate}

\paragraph{考点与思路}
同态保留由加法、乘法和 $0,1$ 写出的等式，却未必保留“非零”“不存在非零零乘积”等条件。零因子的定义涉及非零元素，所以不能只把 $ab=0$ 两边取像就认定像仍为零因子。
\textbf{原文依据：}现有笔记第三章习题8解答：按原六项顺序，分别取像验证；反例依次沿用 $\Z/6\Z\to\Z/3\Z$、$\Z\to\Z/6\Z$、$\Z/6\Z\to\Z/2\Z$、$\Z\to\Z/5\Z$，不更换实例。

\paragraph{所用结论}\toolref{r-homomorphism}、\toolref{r-quotients}。
\begin{proof}[详细证明]
\textbf{(i) 成立。}
若 $a$ 是幂零元，取正整数 $m$ 使 $a^m=0$。由同态保乘法，
\[
\eta(a)^m=\eta(a^m)=\eta(0)=0'.
\]
故 $\eta(a)$ 仍为幂零元。若 $e^2=e$，则
\[
\eta(e)^2=\eta(e^2)=\eta(e),
\]
故 $\eta(e)$ 仍为幂等元。两项结论都不要求满射或保幺。

\textbf{(ii) 不成立。}
考虑映射
\[
\eta:\Z/6\Z\longrightarrow\Z/3\Z,\qquad
\overline{k}_{6}\longmapsto\overline{k}_{3}.
\]
若 $k\equiv\ell\pmod6$，则 $3\mid k-\ell$，所以像不依赖代表元，映射良定义。模运算保留加法和乘法，且 $\eta(\overline1_6)=\overline1_3$；每个 $\overline{k}_3$ 都是 $\overline{k}_6$ 的像，故这是保幺满同态。

在源环中，$\overline4_6\ne0$、$\overline3_6\ne0$，且
\[
\overline4_6\,\overline3_6=\overline{12}_6=\overline0_6,
\]
所以 $\overline4_6$ 是零因子。它的像却是 $\overline1_3$。若 $\overline1_3 y=0$，则 $y=0$，故该像不是零因子。这一反例甚至将零因子映成了非零单位。

\textbf{(iii) 不成立。}
取自然映射
\[
\pi:\Z\twoheadrightarrow\Z/6\Z,\qquad k\longmapsto\overline{k}_6.
\]
该映射保加法、乘法和单位元，并且满射。源环 $\Z$ 是整环，但目标中
\[
\overline2_6\ne0,\qquad \overline3_6\ne0,\qquad
\overline2_6\,\overline3_6=0,
\]
所以 $\Z/6\Z$ 不是整环。

\textbf{(iv) 不成立。}
取映射
\[
\rho:\Z/6\Z\twoheadrightarrow\Z/2\Z,\qquad
\overline{k}_{6}\longmapsto\overline{k}_{2}.
\]
因为 $2\mid6$，同余模 $6$ 蕴含同余模 $2$，映射良定义；模运算说明它保加法、乘法及单位元，每个模 $2$ 的剩余类都有原像，所以满射。

目标 $\Z/2\Z$ 只有两个元素，唯一非零元素是 $\overline1_2$，且 $\overline1_2\,\overline1_2=\overline1_2\ne0$，因此它是整环。源环却有上一项给出的非零零乘积，故不是整环。

\textbf{(v) 成立。}
若 $u$ 是单位，令 $v=u^{-1}$，则 $uv=vu=1$。取像得
\begin{align*}
\eta(u)\eta(v)&=\eta(uv)=\eta(1)=1',\\
\eta(v)\eta(u)&=\eta(vu)=\eta(1)=1'.
\end{align*}
因此 $\eta(u)$ 是单位，且
\[
\eta(u)^{-1}=\eta(u^{-1}).
\]
这一结论只需要保幺同态，不需要满射。

\textbf{(vi) 不成立。}
取自然保幺满同态
\[
\sigma:\Z\twoheadrightarrow\Z/5\Z,\qquad
k\longmapsto\overline{k}_5.
\]
整数 $2$ 不是 $\Z$ 中的单位：若 $2m=1$，则偶数等于奇数，故没有整数解。它的像 $\overline2_5$ 却满足
\[
\overline2_5\,\overline3_5=\overline6_5=\overline1_5,
\]
因而是单位。这说明像可逆不保证原元素可逆。
\end{proof}

\paragraph{复习提示}
六项判断依次为：\textbf{对、错、错、错、对、错}。构造剩余类环之间的映射时，先检验模数之间的整除关系，确保“按原整数取剩余类”的规则良定义。证明零因子的反例时，须同时写明两个因子非零；证明单位的反例时，须给出像的实际逆元。

\clearpage% !TEX root = main.tex
\section{考前速查：结论与证明入口}
下面只列结论和起手步骤，完整论证见各题详解。证明题的“答案”包括理由，不能只写本表中的结论。
\small
\begin{longtable}{p{1.1cm}p{1.0cm}p{12.1cm}}
\toprule 章 & 题号 & 结论与起手步骤\\\midrule\endhead
零&4&依次：全成立；自反与传递；仅传递；对称与传递。\\
零&5&依次：仅对称；自反与对称；全不成立；仅对称。\\
零&6&递归定义与归纳；先证加法，再证分配律引理，最后证乘法。\\
零&7&对$c$归纳证明$a(b+c)=ab+ac$，右分配律由交换律推出。\\
零&8&若$x<y$，写$y=x+t$，$t>0$，则$ay=ax+at>ax$。\\
零&13&共同因子$d$使$a^{2^m}$同时同余$1$与$-1$，故$d\mid2$；取$a=2$构造互素数列。\\
零&14&无平方因子的约数对应不同素因子集合的子集，约数和为$(1-1)^r$。\\
零&23&$x\equiv17\pmod{31}$，逆元$27^{-1}\equiv23\pmod{31}$。\\
零&24&孙子定理原构造：$e_1=231,e_2=330,e_3=210$，$x=2e_1+e_2+5e_3\equiv302\pmod{385}$。\\\midrule
一&1&展开$(ab)^2=a^2b^2$，左右消去得到$ba=ab$。\\
一&2&所有元素等于自身逆元；$ab=(ab)^{-1}=b^{-1}a^{-1}=ba$。\\
一&4&解$ae=a$与$fa=a$；用任意元素的可解表示推广成双侧单位元，再构造逆元。\\
一&6&构造正$n$边形的二面体群$D_n$，阶$2n$，$n>2$时不交换。\\
一&7&共轭保持乘积，$a^{-1}b^ma=(a^{-1}ba)^m$，再对次数归纳。\\
一&8&取逆映射总是双射；保乘法当且仅当群交换。\\
一&9&等价关系当且仅当$S$为子群；等价类为右陪集$Sb$。\\
一&10&$\Z\to n\Z$，$k\mapsto nk$，是加法群同构。\\
一&12&指数$2$时，群外元素的左右陪集都等于$G\setminus H$。\\
一&18&无限阶元素会给出无穷多个不同循环子群；否则群是有限多个有限循环子群的并。\\
一&21&对满射$H\times K\to HK$计数，每个纤维与$H\cap K$一一对应。\\
一&22&正规性给出$MN=NM$；$m\mapsto mN$的核为$M\cap N$。\\
一&29&用$f(a)=1-a$把新运算转为旧运算；新零元$1$，新单位元$0$。\\
一&30&四组例子：$2\Z\subset\Z$；$\{0\}\times\Z\subset2\Z\times\Z$；角块子环；数量矩阵子环。\\
一&32&非零理想含$a\ne0$，再含$a^{-1}a=1$，故为全域。\\
一&34&若$r^n,s^m\in I$，展开$(r-s)^{n+m-1}$，每一项均落在$I$。\\
一&39&令$I=0$使用上题；非交换反例为$E_{12},E_{21}$。\\
一&43&真理想公共零点：紧致性加平方和；零因子：零点集有非空相对开区间。\\
一&44&$|M_n(\F_p)|=p^{n^2}$；$|\GL_n(\F_p)|=\prod_{k=0}^{n-1}(p^n-p^k)$。\\\midrule
二&7&有限个有理数统一分母后，问题化为整数的最大公因数。\\
二&8&有限阶生成元在交换群中能合并同类项，指数取有限范围。\\
二&10&对换生成$S_n$；用共轭把星形对换或相邻对换变成任意对换。\\
二&11&用长轮换共轭$(12)$，得到所有相邻对换。\\
二&12&依教材原题，奇数$n$用$(123),(12\cdots n)$，偶数$n$用$(123),(23\cdots n)$；共轭后逐点加入。\\
二&15&$(H_1\cap H_2)/N=(H_1/N)\cap(H_2/N)$，比较陪集代表元。\\
二&29&固定$g$，$A$与$gB^{-1}$的大小之和超过$|G|$，交集非空，故$g\in AB$。\\
二&30&按不同共轭子群计数，所有共轭共享单位元，覆盖数严格小于$|G|$。\\
二&31&$S_4$有3个西罗2子群和4个西罗3子群；分别列出8阶与3阶子群。\\
二&32&56阶群：若西罗7子群有8个，占48个非单位元；148阶群：西罗37子群唯一。\\
二&36&单位上三角群$U$阶为$p^{n(n-1)/2}$；正规化子为可逆上三角群$B$；个数为$\prod_{i=1}^n(p^i-1)/(p-1)^n$。\\
二&38&共轭同态$G\to\Aut(C_p)$的像阶同时整除$p$幂与$p-1$，故平凡。\\
二&39&$G/Z(G)$循环时，每个元素写成$a^iz$，两个元素直接相乘验证交换。\\
二&40&凯莱嵌入$S_n$，再以符号同态补上两个新点，嵌入$A_{n+2}$。\\
二&41&$[g,n]\in N\cap G^{(1)}=\{e\}$，故$n$与所有$g$交换。\\
二&42&子群的导群逐项包含；商群导群为原导群的像。\\
二&43&先找正规西罗子群；正规部分与商群均可解，使用可解性的扩张引理。\\
二&44&有限$p$群中心非平凡；商去中心后按群阶归纳。\\
二&45&映射$g\mapsto(gH,gK)$的核为$H\cap K$；商群嵌入两个可解群的直积。\\
二&46&$\Inn(G)\cong G/Z(G)$；若内自同构恰2个，则中心商循环，反而推出只有1个。\\
二&47&非交换群取非平凡内自同构；交换群先看取逆映射，指数2时交换两个基向量。\\
二&49&中心平凡；$S_3$自同构作用于三个对换，$S_4$自同构作用于四个西罗3子群的正规化子。\\\midrule
三&1&$(1-ba)^{-1}=1+b(1-ab)^{-1}a$，左右乘积均须验算。\\
三&2&固定$uv=1$，两个右逆之差满足$ut=0$；右逆唯一推出$vu=1$。\\
三&3&令$e=1-vu$，构造$x_k=v+eu^k$；若重复，右乘$v^j$推出$e=0$。\\
三&8&六项判断：对、错、错、错、对、错。反例须检验良定义、满射、保幺与非零条件。\\
\bottomrule
\end{longtable}
\normalsize

\end{document}
